 % 本文件由 scripts/assemble.py 自动装配，勿手改（改 parts/ 后重跑）
% 第8章 宇宙学

% 第8章 Cosmology
\chapter{宇宙学}

\section{最大对称宇宙}

% ---- 书页 323 / p338 ----
当代的宇宙学模型基于这样一种观念：宇宙在各处都大致相同——这一立场有时被称为\textbf{哥白尼原理}（Copernican principle）。乍一看，这种主张似乎很疯狂；比方说，太阳的中心与星际空间那荒凉的寒冷几乎毫无相似之处。但我们把哥白尼原理理解为仅仅适用在最大的尺度上，在那种尺度上密度的局域变化已被平均掉。它在这种尺度上的有效性体现在许多不同的观测之中，例如对星系的计数以及对弥漫 X 射线和 $\gamma$ 射线背景的观测；不过它最清晰的表现是 3K 的宇宙微波背景（cosmic microwave background，CMB）。虽然我们现在知道，微波背景辐射并非完美光滑（其实从来也没有人指望它完美光滑），但对规则性的偏离只有 $10^{-5}$ 的量级或更小，对于大尺度上时空的近似描述来说，这当然是足够好的基础。

哥白尼原理与流形可能具有的两个在数学上更精确的性质相关：各向同性（isotropy）与均匀性（homogeneity）。\textbf{各向同性}作用于流形中某个特定的点，它断言：无论你朝哪个方向看，空间看起来都一样。更正式地说，如果对于 $T_pM$ 中的任意两个矢量 $V$ 和 $W$，都存在 $M$ 的一个等距变换（isometry），使得 $W$ 在此等距变换下的推前与 $V$ 平行（$V$ 本身不被推前），那么流形 $M$ 就称为绕点 $p$ 各向同性。微波背景的观测所指示的，正是空间的各向同性。

\textbf{均匀性}则是度规在整个流形上都相同的陈述。换句话说，给定 $M$ 中任意两点 $p$ 和 $q$，都存在一个把 $p$ 变到 $q$ 的等距变换。注意，均匀性与各向同性之间没有必然的联系；一个流形可以均匀却无一处各向同性（比如通常度规下的 $\mathbf{R}\times S^2$），也可以绕某一点各向同性却不均匀（比如一个锥面，它绕自己的顶点各向同性，但显然不均匀）。另一方面，如果一个空间\emph{处处}各向同性，那么它就是均匀的。同样，如果它绕某一点各向同性并且还均匀，那么它将绕每一点都各向同性。既然各向同性有着充分的观测证据，而哥白尼原理又要我们相信自己并非宇宙的中心、因而别处的观者也应观测到各向同性，那么从今往后我们就同时假定均匀性和各向同性。

% ---- 书页 324 / p339 ----
均匀性和各向同性的用处在于，它们蕴含空间是最大对称的（maximally symmetric）。把各向同性想成转动下的不变性，把均匀性想成经过适当推广的平移下的不变性。于是均匀性和各向同性合在一起就蕴含：一个空间具有其可能的最大数目的 Killing 矢量。哥白尼原理的一个极端应用，是坚持认为时空本身就是最大对称的。事实上结果将表明并非如此；从观测上我们知道，宇宙在\emph{空间}上是均匀且各向同性的，但在全部\emph{时空}上却不是。不过，从考察最大对称的时空出发是有趣的（它们毕竟只是“仅空间最大对称”这一更普遍情形的特例）。我们将会看到，在某种意义上，这样的宇宙是广义相对论的“基态”。与本章后面的部分相比，这一讨论与观测到的宇宙关系较远；注重经验的读者不妨直接跳到下一节去。

我们在第 3 章中曾提到，带度规 $g_{\mu\nu}$ 的最大对称 $n$ 维流形，其 Riemann 张量可以写成
\begin{equation*}
R_{\rho\sigma\mu\nu} = \kappa\left(g_{\rho\mu}g_{\sigma\nu} - g_{\rho\nu}g_{\sigma\mu}\right),
\tag{8.1}
\end{equation*}
其中 $\kappa$ 是 Ricci 曲率的归一化量度，
\begin{equation*}
\kappa = \frac{R}{n(n-1)},
\tag{8.2}
\end{equation*}
而 Ricci 标量 $R$ 在流形上将是一个常数。由于在任何一个单点上我们总可以把度规化为规范形式（$g_{\mu\nu}=\eta_{\mu\nu}$），最大对称流形的种类在局域上就由度规的号差以及常数 $\kappa$ 的符号来刻画。需要“局域”这个限定词，是为了顾及可能的全局差异，比如平面与环面之间的差异。我们感兴趣的是号差为 $(-+++)$ 的度规。对于曲率为零（$\kappa=0$）的情形，最大对称时空是众所周知的；它恰恰就是 Minkowski 空间，其度规为
\begin{equation*}
\mathrm{d}s^2 = -\mathrm{d}t^2 + \mathrm{d}x^2 + \mathrm{d}y^2 + \mathrm{d}z^2.
\tag{8.3}
\end{equation*}
Minkowski 空间的共形图在附录 H 中导出。

曲率为正（$\kappa>0$）的最大对称时空称为 \textbf{de Sitter 空间}（de Sitter space）。考虑一个五维 Minkowski 空间，度规为 $\mathrm{d}s_5^2 = -\mathrm{d}u^2 + \mathrm{d}x^2 + \mathrm{d}y^2 + \mathrm{d}z^2 + \mathrm{d}w^2$，并在其中嵌入一个由下式给出的双曲面（hyperboloid）
\begin{equation*}
-u^2 + x^2 + y^2 + z^2 + w^2 = \alpha^2.
\tag{8.4}
\end{equation*}
现在按如下方式在双曲面上诱导坐标 $\{t,\chi,\theta,\phi\}$：
\begin{equation*}
\begin{aligned}
u &= \alpha\sinh(t/\alpha)\\
w &= \alpha\cosh(t/\alpha)\cos\chi\\
x &= \alpha\cosh(t/\alpha)\sin\chi\cos\theta\\
y &= \alpha\cosh(t/\alpha)\sin\chi\sin\theta\cos\phi\\
z &= \alpha\cosh(t/\alpha)\sin\chi\sin\theta\sin\phi.
\end{aligned}
\tag{8.5}
\end{equation*}
% ---- 书页 325 / p340 ----（原书 324/325 分页恰落在本式内：(u,w) 在书页 324，(x,y,z) 及其后文字在书页 325；两块已按原书合并为单个五行公式） ----
于是双曲面上的度规为
\begin{equation*}
\mathrm{d}s^2 = -\mathrm{d}t^2 + \alpha^2\cosh^2(t/\alpha)\left[\mathrm{d}\chi^2 + \sin^2\chi\left(\mathrm{d}\theta^2 + \sin^2\theta\,\mathrm{d}\phi^2\right)\right].
\tag{8.6}
\end{equation*}
我们认出圆括号中的表达式是二维球面上的度规 $\mathrm{d}\Omega_2^2$，而方括号中的表达式是三维球面上的度规 $\mathrm{d}\Omega_3^2$。这样，de Sitter 空间描述的是一个先是收缩、在 $t=0$ 时达到最小尺寸、然后再重新膨胀的空间三维球面。当然，这一特殊描述承袭自某个特定的坐标系；我们将看到，同样有效的其他描述也是存在的。

这些坐标覆盖整个流形。一般你可以这样检验这一点：比如追踪测地线在坐标系边缘附近的行为；如果坐标不完备，测地线看起来就会在有限的仿射参量处终止。因此 de Sitter 空间的拓扑是 $\mathbf{R}\times S^3$。这使得导出共形图变得非常简单，因为构造共形图的关键一步，就是把度规写成它与 Einstein 静态宇宙（一个拓扑为 $\mathbf{R}\times S^3$、描述一个随时间推移保持半径不变的空间三维球面的时空）共形相关联的形式。考虑从 $t$ 到 $t'$ 的坐标变换
\begin{equation*}
\cosh(t/\alpha) = \frac{1}{\cos(t')}.
\tag{8.7}
\end{equation*}
度规 (8.6) 于是变为
\begin{equation*}
\mathrm{d}s^2 = \frac{\alpha^2}{\cos^2(t')}\,\mathrm{d}\bar{s}^2,
\tag{8.8}
\end{equation*}
其中 $\mathrm{d}\bar{s}^2$ 表示 Einstein 静态宇宙上的度规，
\begin{equation*}
\mathrm{d}\bar{s}^2 = -(\mathrm{d}t')^2 + \mathrm{d}\chi^2 + \sin^2\chi\,\mathrm{d}\Omega_2^2.
\tag{8.9}
\end{equation*}
新时间坐标的范围是
\begin{equation*}
-\pi/2 < t' < \pi/2.
\tag{8.10}
\end{equation*}
de Sitter 空间的共形图就是与 de Sitter 共形相关的那一片 Einstein 静态宇宙区域的表示。它看起来像一个正方形，如图 8.1 所示。常数 $t'$ 的类空切片代表一个三维球面；左右两条边缘处的虚线是这个球面的北极和南极。对角线代表类光射线；在过去无限远释放的一个光子，恰好会在未来无限远时到达球面上正对面的那一点（对跖点）。

\begin{figure}[H]
\centering
\includegraphics[width=0.72\textwidth]{fig_8.1.pdf}
\caption*{图 8.1\quad de Sitter 时空的共形图。类空切片是三维球面，因此图上的点代表二维球面，只有左右两条边缘上的点除外——它们代表单个点。}
\end{figure}

% ---- 书页 326 / p341 ----
请记住，时空之所以向过去和未来“终结”，仅仅是共形变换的魔力所致；真实的 de Sitter 空间向未来和过去无限延伸。还要注意，两个点可以拥有完全断开的未来（或过去）光锥；这反映了这样一个事实：球状空间截面膨胀得如此之快，以至于来自一点的光永远无法与来自另一点的光相接触。

一个类似的双曲面构造揭示了最大对称而 $\kappa<0$ 的时空，即所谓的 \textbf{反 de Sitter 空间}（anti-de Sitter space）。从一个虚构的五维平直流形出发，其度规为 $\mathrm{d}s_5^2 = -\mathrm{d}u^2 - \mathrm{d}v^2 + \mathrm{d}x^2 + \mathrm{d}y^2 + \mathrm{d}z^2$，并嵌入一个由下式给出的双曲面
\begin{equation*}
-u^2 - v^2 + x^2 + y^2 + z^2 = -\alpha^2.
\tag{8.11}
\end{equation*}
注意这里所有的负号。然后我们可以在双曲面上诱导坐标 $\{t',\rho,\theta,\phi\}$：
\begin{equation*}
\begin{aligned}
u &= \alpha\sin(t')\cosh(\rho)\\
v &= \alpha\cos(t')\cosh(\rho)\\
x &= \alpha\sinh(\rho)\cos\theta\\
y &= \alpha\sinh(\rho)\sin\theta\cos\phi\\
z &= \alpha\sinh(\rho)\sin\theta\sin\phi,
\end{aligned}
\tag{8.12}
\end{equation*}
由此得到这个双曲面上如下形式的度规：
\begin{equation*}
\mathrm{d}s^2 = \alpha^2\left(-\cosh^2(\rho)\,\mathrm{d}t'^2 + \mathrm{d}\rho^2 + \sinh^2(\rho)\,\mathrm{d}\Omega_2^2\right).
\tag{8.13}
\end{equation*}
这些坐标有一个奇怪的特点，即 $t'$ 是周期的。由 (8.12)，$t'$ 与 $t'+2\pi$ 代表双曲面上的同一个位置。由于 $\partial_{t'}$ 处处类时，当 $t'$ 增加时保持 $\{\rho,\theta,\phi\}$ 不变的曲线将是一条闭合类时曲线。然而，这并不是时空的内禀性质，只是我们从特定嵌入导出度规的方式所带来的假象。我们完全可以转而考虑这个流形的“覆盖空间”（covering space），即度规由 (8.13) 给出、但允许 $t'$ 从 $-\infty$ 变到 $\infty$ 的时空。这个空间中没有闭合类时曲线，我们就把这一点取作反 de Sitter 空间的定义。

为了导出共形图，进行一个与 de Sitter 情形类似的坐标变换，不过这次作用在径向坐标上：
\begin{equation*}
\cosh(\rho) = \frac{1}{\cos\chi},
\tag{8.14}
\end{equation*}
于是
\begin{equation*}
\mathrm{d}s^2 = \frac{\alpha^2}{\cos^2\chi}\,\mathrm{d}\bar{s}^2,
\tag{8.15}
\end{equation*}

% ---- 书页 327 / p342 ----
\begin{figure}[H]
\centering
\includegraphics[width=0.72\textwidth]{fig_8.2.pdf}
\caption*{图 8.2\quad 反 de Sitter 时空的共形图。类空切片具有 $\mathbf{R}^3$ 的拓扑，我们用极坐标把它表示出来，因此图上的点代表二维球面，只有左边的点除外——它们代表空间原点处的单个点。无限远是右侧的一个类时曲面。}
\end{figure}

其中 $\mathrm{d}\bar{s}^2$ 表示 Einstein 静态宇宙的度规 (8.9)。与 de Sitter 不同，径向坐标现在出现在共形因子之中。此外，对反 de Sitter 而言，$t'$ 坐标从负无穷一直到正无穷，而径向坐标的范围是
\begin{equation*}
0 \leq \chi < \frac{\pi}{2}.
\tag{8.16}
\end{equation*}
这样，反 de Sitter 空间与 Einstein 静态宇宙的一半共形相关联。其共形图示于图 8.2，图中画出了几条通过点 $t'=0$、$\chi=0$ 的有代表性的类时和类空测地线。由于 $\chi$ 只到 $\pi/2$ 而不是一直到 $\pi$，这个时空的类空切片具有 $S^3$ 中一个半球面内部的拓扑；也就是说，它在拓扑上是 $\mathbf{R}^3$（因而整个时空具有拓扑 $\mathbf{R}^4$）。注意，我们用极坐标画出这张图：左边的一点代表空间原点处的一个点，而右边的一点代表空间无限远处的一个二维球面。另一种流行的表示方法是画出时空的横截面，让空间原点位于正中间，而右边和左边合起来构成空间无限远。

反 de Sitter 的一个有趣特点是，无限远以一个类时超曲面的形式出现，由 $\chi=\pi/2$ 定义。由于无限远是类时的，这个空间不是全局双曲的，我们就没有一个以类空切片上给定的信息来表述的适定初值问题，因为信息总是可以“从无穷远流入”。另一个有趣的特点是，指数映射并不映满整个时空；如图中所画的那样，从指定点出发的测地线并不覆盖整个流形。指向未来的类时测地线，如图所示，起初可以从 $t'=0$、$\chi=0$ 沿径向向外运动，但最终会重新聚焦到点 $t'=\pi$、$\chi=0$，然后再次沿径向向外运动。

% ---- 书页 328 / p343 ----
顺便说一句，我们忍不住要指出：无限远的类时性质促成了弦论的一个引人注目的特征——“AdS/CFT 对应”（AdS/CFT correspondence）。这里 AdS 当然就是我们一直在讨论的反 de Sitter 空间，而 CFT 指的是定义在边界上的共形不变的场论［对 $n$ 维 AdS 来说，这个边界本身就是一个 $(n-1)$ 维时空］。AdS/CFT 对应表明，在某个极限下，AdS 背景上的量子引力（或其超对称版本）与定义在边界上的共形不变的非引力场论之间存在等价关系。关于非引力的量子场论我们知道得很多，而对量子引力却知之甚少，因此这一对应（如果它是正确的——这看起来很有可能，但尚未被证明）将揭示量子引力中可能发生之事的大量信息。\footnote{全面的综述文章见 O.~Aharony, S.S.~Gubser, J.M.~Maldacena, H.~Ooguri, and Y.~Oz, \emph{Phys.\ Rept.} \textbf{323}, 183 (2000), \texttt{http://arxiv.org/hep-th/9905111}。}

于是我们有了三个最大对称的时空：Minkowski（$\kappa=0$）、de Sitter（$\kappa>0$）和反 de Sitter（$\kappa<0$）。它们中有哪一个是真实世界的有用模型吗？进而言之，它们是 Einstein 方程的解吗？先对 (8.1) 给出的 Riemann 张量取迹，并限定到四维：
\begin{equation*}
R_{\mu\nu} = 3\kappa g_{\mu\nu},\qquad R = 12\kappa.
\tag{8.17}
\end{equation*}
可见在最大对称空间中，Ricci 张量与度规成正比。具有这种性质的时空有时被称为 Einstein 空间（Einstein space）；而 Einstein 静态宇宙并\emph{不}是 Einstein 空间的一个例子，这一点有时会令人混淆。更糟的是，我们后面还会遇到 Einstein–de Sitter 宇宙学，它既与 Einstein 空间无关，也与 Einstein 静态宇宙或 de Sitter 空间无关。Einstein 张量为
\begin{equation*}
G_{\mu\nu} = R_{\mu\nu} - \frac{1}{2}Rg_{\mu\nu} = -3\kappa g_{\mu\nu}.
\tag{8.18}
\end{equation*}
因此，Einstein 方程 $G_{\mu\nu}=8\pi G T_{\mu\nu}$ 蕴含（在最大对称时空中如此，并非普遍情形）能量-动量张量与度规成正比：
\begin{equation*}
T_{\mu\nu} = -\frac{3\kappa}{8\pi G}\,g_{\mu\nu}.
\tag{8.19}
\end{equation*}
这样的能量-动量张量对应于真空能或宇宙学常数，如第 4 章所讨论的。能量密度和压强为
\begin{equation*}
\rho = -p = \frac{3\kappa}{8\pi G}.
\tag{8.20}
\end{equation*}
若 $\rho$ 为正，我们得到 de Sitter 解；若 $\rho$ 为负，我们得到反 de Sitter。但在我们的宇宙中，除了可能有真空能之外，还有普通的物质和辐射。我们的最大对称时空并不相容于
% ---- 书页 329 / p344 ----
动力学上可观的物质和/或辐射。此外，由于我们观测到宇宙中的可见物质正在彼此分开（宇宙正在膨胀，下面会讨论），物质的密度在过去更高；所以即使物质对总能量的贡献在今天可以忽略，它在更早的宇宙中也曾经是可观的。因此，最大对称时空并不是真实世界的合理模型。不过，它们确实代表了在没有任何普通物质或引力辐射时 Einstein 方程的（局域）唯一解；正是在这个意义上，它们可以被看作广义相对论的基态。

\section{Robertson–Walker 度规}

为了描述真实世界，我们不得不放弃“完美的”哥白尼原理——它蕴含空间和时间处处都对称——而假定某种更宽容的东西。事实证明，直接假定宇宙在\emph{空间上}均匀且各向同性、但随时间演化，既是直截了当的，又与观测一致。在广义相对论中，这转译为如下陈述：宇宙可以叶状分层（foliated）为一些类空切片，使得每个三维切片都是最大对称的。于是我们把时空取为 $\mathbf{R}\times\Sigma$，其中 $\mathbf{R}$ 代表时间方向，而 $\Sigma$ 是一个最大对称的三维流形。时空度规因此取如下形式
\begin{equation*}
\mathrm{d}s^2 = -\mathrm{d}t^2 + R^2(t)\,\mathrm{d}\sigma^2,
\tag{8.21}
\end{equation*}
其中 $t$ 是类时坐标，$R(t)$ 是一个称为\textbf{尺度因子}（scale factor）的函数，而 $\mathrm{d}\sigma^2$ 是 $\Sigma$ 上的度规，它可以表示为
\begin{equation*}
\mathrm{d}\sigma^2 = \gamma_{ij}(u)\,\mathrm{d}u^i\,\mathrm{d}u^j,
\tag{8.22}
\end{equation*}
其中 $(u^1,u^2,u^3)$ 是 $\Sigma$ 上的坐标，$\gamma_{ij}$ 是一个最大对称的三维度规。尺度因子告诉我们类空切片 $\Sigma$ 在时刻 $t$ 有多大。（不要把它与标量曲率混淆。）这里使用的坐标，其度规没有交叉项 $\mathrm{d}t\,\mathrm{d}u^i$，且 $\mathrm{d}t^2$ 的系数与 $u^i$ 无关，称为\textbf{共动坐标}（comoving coordinates），它是附录 D 讨论的高斯法坐标的一个特例。保持在常数 $u^i$ 处的观者也称为“共动的”。只有共动观者才会认为宇宙看起来是各向同性的；事实上在地球上我们并不完全共动，结果便是我们在宇宙微波背景中看到偶极各向异性，它来源于常规的多普勒效应。

因此我们感兴趣的是最大对称的欧几里得三维度规 $\gamma_{ij}$。我们知道，最大对称的度规满足
\begin{equation*}
{}^{(3)}R_{ijkl} = k\left(\gamma_{ik}\gamma_{jl} - \gamma_{il}\gamma_{jk}\right),
\tag{8.23}
\end{equation*}
其中为了今后的方便我们引入了
\begin{equation*}
k = {}^{(3)}R/6,
\tag{8.24}
\end{equation*}

% ---- 书页 330 / p345 ----
并且在 Riemann 张量上放一个上标 $(3)$，提醒我们它与三维度规 $\gamma_{ij}$ 相联系，而不是与整个时空的度规相联系。Ricci 张量便是
\begin{equation*}
{}^{(3)}R_{jl} = 2k\gamma_{jl}.
\tag{8.25}
\end{equation*}
如果空间要是最大对称的，那它肯定是球对称的。从对 Schwarzschild 解的探索中，我们已经对球对称空间有所了解；度规可以写成
\begin{equation*}
\mathrm{d}\sigma^2 = \gamma_{ij}\,\mathrm{d}u^i\,\mathrm{d}u^j = e^{2\beta(\bar{r})}\,\mathrm{d}\bar{r}^2 + \bar{r}^2\,\mathrm{d}\Omega^2,
\tag{8.26}
\end{equation*}
其中 $\bar{r}$ 是径向坐标，二维球面上的度规一如既往是 $\mathrm{d}\Omega^2 = \mathrm{d}\theta^2 + \sin^2\theta\,\mathrm{d}\phi^2$。这种度规的 Ricci 张量分量可以从 (5.14)——静态球对称时空的 Ricci 张量——通过令 $\alpha=0$ 和 $r=\bar{r}$ 得到，即
\begin{equation*}
\begin{aligned}
{}^{(3)}R_{11} &= \frac{2}{\bar{r}}\,\partial_1\beta\\
{}^{(3)}R_{22} &= e^{-2\beta}\left(\bar{r}\,\partial_1\beta - 1\right) + 1\\
{}^{(3)}R_{33} &= \left[e^{-2\beta}\left(\bar{r}\,\partial_1\beta - 1\right) + 1\right]\sin^2\theta.
\end{aligned}
\tag{8.27}
\end{equation*}
我们利用 (8.25) 令它们与度规成正比，就可以解出 $\beta(\bar{r})$：
\begin{equation*}
\beta = -\frac{1}{2}\ln\left(1-k\bar{r}^2\right),
\tag{8.28}
\end{equation*}
由此得到三维曲面 $\Sigma$ 上的度规
\begin{equation*}
\mathrm{d}\sigma^2 = \frac{\mathrm{d}\bar{r}^2}{1-k\bar{r}^2} + \bar{r}^2\,\mathrm{d}\Omega^2.
\tag{8.29}
\end{equation*}
从 (8.24) 可以注意到，$k$ 的值决定了空间曲面的曲率，从而决定了空间曲面的大小。通常把它归一化，使得
\begin{equation*}
k \in \{+1, 0, -1\},
\tag{8.30}
\end{equation*}
而把流形的物理尺寸吸收进尺度因子 $R(t)$。

$k=-1$ 的情形对应于 $\Sigma$ 上曲率为负的常数，有时称为\textbf{开放}（open）；$k=0$ 的情形对应于 $\Sigma$ 上没有曲率，称为\textbf{平直}（flat）；$k=+1$ 的情形对应于 $\Sigma$ 上曲率为正，有时称为\textbf{闭合}（closed）。引入一个由下式定义的新径向坐标 $\chi$，得到度规的另一种形式，可以使这些情形的物理解释更加清楚：
\begin{equation*}
\mathrm{d}\chi = \frac{\mathrm{d}\bar{r}}{\sqrt{1-k\bar{r}^2}}.
\tag{8.31}
\end{equation*}

% ---- 书页 331 / p346 ----
这可以积分得到
\begin{equation*}
\bar{r} = S_k(\chi),
\tag{8.32}
\end{equation*}
其中
\begin{equation*}
S_k(\chi) \equiv
\begin{cases}
\sin(\chi), & k = +1\\
\chi, & k = 0\\
\sinh(\chi), & k = -1,
\end{cases}
\tag{8.33}
\end{equation*}
于是
\begin{equation*}
\mathrm{d}\sigma^2 = \mathrm{d}\chi^2 + {S_k}^2(\chi)\,\mathrm{d}\Omega^2.
\tag{8.34}
\end{equation*}
对平直情形 $k=0$，$\Sigma$ 上的度规变为
\begin{equation*}
\begin{aligned}
\mathrm{d}\sigma^2 &= \mathrm{d}\chi^2 + \chi^2\,\mathrm{d}\Omega^2\\
&= \mathrm{d}x^2 + \mathrm{d}y^2 + \mathrm{d}z^2,
\end{aligned}
\tag{8.35}
\end{equation*}
这就是平直的欧几里得空间。整体上，它可以描述 $\mathbf{R}^3$ 或者更复杂的流形，比如三维环面 $S^1\times S^1\times S^1$。对闭合情形 $k=+1$，我们有
\begin{equation*}
\mathrm{d}\sigma^2 = \mathrm{d}\chi^2 + \sin^2\chi\,\mathrm{d}\Omega^2,
\tag{8.36}
\end{equation*}
这是三维球面的度规。在这种情形下，唯一可能的整体结构是完全的三维球面（不可定向流形 $\mathbf{RP}^3$ 除外，它由认同 $S^3$ 上的对跖点得到）。最后，对开放的 $k=-1$ 情形，我们得到
\begin{equation*}
\mathrm{d}\sigma^2 = \mathrm{d}\chi^2 + \sinh^2\chi\,\mathrm{d}\Omega^2.
\tag{8.37}
\end{equation*}
这是常负曲率三维空间的度规，是 3.9 节讨论的那个双曲面的推广。整体上，这样的空间可以无限延伸（这正是“开放”一词的由来），但它也可以描述一个非单连通的紧致空间（所以“开放”实在算不上最准确的描述）。

时空的度规描述的正是这些最大对称超曲面之一在尺寸上的演化，它可以写成
\begin{equation*}
\mathrm{d}s^2 = -\mathrm{d}t^2 + R^2(t)\left[\frac{\mathrm{d}\bar{r}^2}{1-k\bar{r}^2} + \bar{r}^2\,\mathrm{d}\Omega^2\right].
\tag{8.38}
\end{equation*}
这就是 \textbf{Robertson–Walker（RW）度规}。我们还没有用到 Einstein 方程；它将决定尺度因子 $R(t)$ 的行为。注意，代换

% ---- 书页 332 / p347 ----
\begin{equation*}
\begin{aligned}
R &\to \lambda^{-1}R\\
\bar{r} &\to \lambda\bar{r}\\
k &\to \lambda^{-2}k
\end{aligned}
\tag{8.39}
\end{equation*}
使 (8.38) 保持不变。因此我们可以选取方便的归一化。在曲率 $k$ 归一化为 $\{+1,0,-1\}$ 的那些变量中，尺度因子具有距离的量纲，而径向坐标 $\bar{r}$（或 $\chi$）实际上是无量纲的；这是最流行的选择。我们偏要反其道而行之，改用无量纲的尺度因子
\begin{equation*}
a(t) = \frac{R(t)}{R_0},
\tag{8.40}
\end{equation*}
一个具有距离量纲的坐标
\begin{equation*}
r = R_0\bar{r},
\tag{8.41}
\end{equation*}
以及一个具有 (长度)$^{-2}$ 量纲的曲率参数
\begin{equation*}
\kappa = \frac{k}{R_0^2}.
\tag{8.42}
\end{equation*}
注意，$\kappa$ 可以取任意值，而不仅仅是 $\{+1,0,-1\}$。在这些变量中，Robertson–Walker 度规为
\begin{equation*}
\boxed{\;\mathrm{d}s^2 = -\mathrm{d}t^2 + a^2(t)\left[\frac{\mathrm{d}r^2}{1-\kappa r^2} + r^2\,\mathrm{d}\Omega^2\right].\;}
\tag{8.43}
\end{equation*}
要转换到更常见的记号，只需代入关系式 (8.40)、(8.41) 和 (8.42)。

有了度规在手，我们就可以着手计算联络系数和曲率张量。令 $\dot{a}\equiv \mathrm{d}a/\mathrm{d}t$，Christoffel 符号为
\begin{equation*}
\begin{array}{ll}
\displaystyle \Gamma^0_{11} = \frac{a\dot{a}}{1-\kappa r^2} & \displaystyle \Gamma^1_{11} = \frac{\kappa r}{1-\kappa r^2}\\[10pt]
\displaystyle \Gamma^0_{22} = a\dot{a}r^2 & \displaystyle \Gamma^0_{33} = a\dot{a}r^2\sin^2\theta\\[6pt]
\displaystyle \Gamma^1_{01} = \Gamma^2_{02} & \displaystyle \Gamma^3_{03} = \frac{\dot{a}}{a}\\[6pt]
\displaystyle \Gamma^1_{22} = -r\left(1-\kappa r^2\right) & \displaystyle \Gamma^1_{33} = -r\left(1-\kappa r^2\right)\sin^2\theta\\[6pt]
\displaystyle \Gamma^2_{12} = \Gamma^3_{13} = \frac{1}{r} & \\[6pt]
\displaystyle \Gamma^2_{33} = -\sin\theta\cos\theta & \displaystyle \Gamma^3_{23} = \cot\theta,
\end{array}
\tag{8.44}
\end{equation*}

% ---- 书页 333 / p348 ----
凡未明确写出的，或者为零，或者由对称性与这些相联系。Ricci 张量的非零分量为
\begin{equation*}
\begin{aligned}
R_{00} &= -3\frac{\ddot{a}}{a}\\[4pt]
R_{11} &= \frac{a\ddot{a} + 2\dot{a}^2 + 2\kappa}{1-\kappa r^2}\\[4pt]
R_{22} &= r^2\left(a\ddot{a} + 2\dot{a}^2 + 2\kappa\right)\\
R_{33} &= r^2\left(a\ddot{a} + 2\dot{a}^2 + 2\kappa\right)\sin^2\theta,
\end{aligned}
\tag{8.45}
\end{equation*}
而 Ricci 标量则为
\begin{equation*}
R = 6\left[\frac{\ddot{a}}{a} + \left(\frac{\dot{a}}{a}\right)^2 + \frac{\kappa}{a^2}\right].
\tag{8.46}
\end{equation*}

\section{弗里德曼方程}

RW 度规对尺度因子 $a(t)$ 的任何行为都有定义；我们的下一步将是把它代入 Einstein 方程，导出把尺度因子与宇宙的能量-动量联系起来的弗里德曼方程（组）（Friedmann equation(s)）。我们将选择用理想流体来模型化物质和能量。显然，如果某种流体在某一个参考系中是各向同性的，而它导致的度规在某一个参考系中是各向同性的，那么这两个参考系将会重合；也就是说，流体在共动坐标中静止。四速度就是
\begin{equation*}
U^\mu = (1, 0, 0, 0),
\tag{8.47}
\end{equation*}
而能量-动量张量
\begin{equation*}
T_{\mu\nu} = (\rho + p)U_\mu U_\nu + pg_{\mu\nu}
\tag{8.48}
\end{equation*}
变为
\begin{equation*}
T_{\mu\nu} = \left(\begin{array}{cccc}
\rho & 0 & 0 & 0\\
0 & & & \\
0 & & g_{ij}p & \\
0 & & &
\end{array}\right).
\tag{8.49}
\end{equation*}
升起一个指标后，它取如下方便的形式：
\begin{equation*}
T^\mu{}_\nu = \mathrm{diag}(-\rho, p, p, p).
\tag{8.50}
\end{equation*}
注意，迹由下式给出：
\begin{equation*}
T = T^\mu{}_\mu = -\rho + 3p.
\tag{8.51}
\end{equation*}
在代入 Einstein 方程之前，先考察能量守恒方程的零分量是有教益的：

% ---- 书页 334 / p349 ----
\begin{equation*}
\begin{aligned}
0 &= \nabla_\mu T^\mu{}_{0}\\
&= \partial_\mu T^\mu{}_{0} + \Gamma^\mu_{\mu\lambda}T^\lambda{}_{0} - \Gamma^\lambda_{\mu 0}T^\mu{}_{\lambda}\\
&= -\partial_0\rho - 3\frac{\dot{a}}{a}\left(\rho + p\right).
\end{aligned}
\tag{8.52}
\end{equation*}
为了取得进展，我们可以选取一个\textbf{物态方程}（equation of state），即 $\rho$ 与 $p$ 之间的一个关系。与宇宙学相关的理想流体常常服从简单的物态方程
\begin{equation*}
p = w\rho,
\tag{8.53}
\end{equation*}
其中 $w$ 是与时间无关的常数。当然，无论它是否保持为常数，我们都可以随意定义参数 $w=p/\rho$；然而，如果 $w$ 在变化，把 $p=w\rho$ 称为“物态方程”就不那么名副其实了。能量守恒方程变为
\begin{equation*}
\boxed{\;\frac{\dot{\rho}}{\rho} = -3(1+w)\frac{\dot{a}}{a}.\;}
\tag{8.54}
\end{equation*}
如果 $w$ 是常数，这可以积分得到
\begin{equation*}
\rho \propto a^{-3(1+w)}.
\tag{8.55}
\end{equation*}
要想了解 $w$ 允许取什么样的值，请参考第 4 章关于能量条件的讨论。类光主能量条件（Null Dominant Energy Condition）允许任何符号的真空能，但除此之外要求物质不能使真空失稳，它蕴含
\begin{equation*}
|w| \leq 1.
\tag{8.56}
\end{equation*}
这个要求绝非不可动摇的金科玉律，但作为研究真实世界中可能发生之事的出发点，它看起来是明智而保守的。

宇宙学流体最流行的两个例子称为\textbf{物质}（matter）和\textbf{辐射}（radiation）。物质是任何一组无碰撞、非相对论性的粒子，其压强基本上为零：
\begin{equation*}
p_{\mathrm{M}} = 0.
\tag{8.57}
\end{equation*}
例子包括普通的恒星和星系，对它们来说，压强与能量密度相比可以忽略。物质也称为尘埃（dust），能量密度主要由物质构成的宇宙称为\textbf{物质主导}的（matter-dominated）。物质中的能量密度按
\begin{equation*}
\rho_{\mathrm{M}} \propto a^{-3}.
\tag{8.58}
\end{equation*}
衰减。这可以直接解释为：随着宇宙膨胀，粒子数密度在减少。对物质而言，能量密度由静能主导，
% ---- 书页 335 / p350 ----
而静能与数密度成正比。辐射既可以描述真实的电磁辐射，也可以描述相对速度足够接近光速、以至于（至少就其物态方程而言）与光子无法区分的有质量粒子。虽然各向同性的相对论性粒子气体是一种理想流体，因而具有由 (8.48) 给出的能量-动量张量，我们也知道，电磁学的 $T_{\mu\nu}$ 可以用场强表达为
\begin{equation*}
T^{\mu\nu} = F^{\mu\lambda}F^{\nu}{}_{\lambda} - \frac{1}{4}g^{\mu\nu}F^{\lambda\sigma}F_{\lambda\sigma}.
\tag{8.59}
\end{equation*}
它的迹由下式给出：
\begin{equation*}
T^\mu{}_\mu = F^{\mu\lambda}F_{\mu\lambda} - \frac{1}{4}(4)F^{\lambda\sigma}F_{\lambda\sigma} = 0.
\tag{8.60}
\end{equation*}
但它也必须等于 (8.51)，所以物态方程为
\begin{equation*}
p_{\mathrm{R}} = \frac{1}{3}\rho_{\mathrm{R}}.
\tag{8.61}
\end{equation*}
能量密度大部分以辐射形式存在的宇宙称为\textbf{辐射主导}的（radiation-dominated）。辐射中的能量密度按
\begin{equation*}
\rho_{\mathrm{R}} \propto a^{-4}.
\tag{8.62}
\end{equation*}
衰减。这样，辐射中的能量密度比物质中的衰减得稍快一些；这是因为光子数密度的减少方式与非相对论性粒子数密度的减少方式相同，但单个光子在红移时还会以 $a^{-1}$ 的方式损失能量，这一点我们稍后会看到。同样，有质量但相对论性的粒子在共动坐标中“慢下来”时也会损失能量。我们相信，今天辐射的能量密度远小于物质的能量密度，$\rho_{\mathrm{M}}/\rho_{\mathrm{R}}\sim 10^{3}$。然而在过去，宇宙要小得多，在极早的时刻，辐射的能量密度本会占据主导。

正如我们讨论过的，真空能也取理想流体的形式，其物态方程为 $p_{\Lambda}=-\rho_{\Lambda}$。其能量密度为常数，
\begin{equation*}
\rho_{\Lambda} \propto a^{0}.
\tag{8.63}
\end{equation*}
由于物质和辐射的能量密度随着宇宙膨胀而减少，如果存在非零的真空能，那么只要宇宙不开始收缩，从长远来看它往往会占上风。如果发生这种情况，我们就说宇宙变得\textbf{真空主导}（vacuum-dominated）。de Sitter 空间和反 de Sitter 空间都是真空主导的解。

现在我们转向 Einstein 方程。回想它可以写成 (4.45) 的形式：
\begin{equation*}
R_{\mu\nu} = 8\pi G\left(T_{\mu\nu} - \frac{1}{2}g_{\mu\nu}T\right).
\tag{8.64}
\end{equation*}
% ==== tail：8.3 节（The Friedmann Equation）收尾，承接书页 335 末的 (8.64)，
% ==== 至书页 337 末止；置于本文件第一个 \section 之前 ====
% ---- 书页 336 / p351 ----
$\mu\nu=00$ 分量方程为
\begin{equation*}
-3\frac{\ddot{a}}{a} = 4\pi G(\rho + 3p),
\tag{8.65}
\end{equation*}
而 $\mu\nu=ij$ 各分量方程给出
\begin{equation*}
\frac{\ddot{a}}{a} + 2\left(\frac{\dot{a}}{a}\right)^{2} + 2\frac{\kappa}{a^{2}} = 4\pi G(\rho - p).
\tag{8.66}
\end{equation*}
由于各向同性，由 $\mu\nu=ij$ 只能得到一个独立的方程。我们可以用式 (8.65) 消去式 (8.66) 中的二阶导数，再稍加整理，得到
\begin{equation*}
\boxed{\,\left(\frac{\dot{a}}{a}\right)^{2} = \frac{8\pi G}{3}\rho - \frac{\kappa}{a^{2}},\,}
\tag{8.67}
\end{equation*}
以及
\begin{equation*}
\boxed{\,\frac{\ddot{a}}{a} = -\frac{4\pi G}{3}(\rho + 3p).\,}
\tag{8.68}
\end{equation*}
这两条方程合称\textbf{弗里德曼方程}（Friedmann equations），而具有式 (8.43) 的形式且满足这些方程的度规，定义了 Friedmann--Robertson--Walker（FRW）宇宙。事实上，如果我们知道 $\rho$ 对 $a$ 的依赖关系，其中第一条式 (8.67) 就足以解出 $a(t)$；当你听到人们提起\emph{那条}弗里德曼方程时，指的就是它，而式 (8.68) 有时被称为\emph{第二}弗里德曼方程。

有一大堆与宇宙学参数相联系的术语，我们在这里只介绍最基本的部分。膨胀的速率由\textbf{Hubble 参量}（Hubble parameter）表征，
\begin{equation*}
\boxed{\,H = \frac{\dot{a}}{a}.\,}
\tag{8.69}
\end{equation*}

% ---- 书页 337 / p352 ----
Hubble 参量在当前时期的值就是 Hubble 常数，$H_0$。目前的测量使我们相信，Hubble 常数为 $70\pm10$ km/sec/Mpc。（Mpc 代表兆秒差距（megaparsec），即 $3.09\times10^{24}$ cm。）由于这个数值仍存在一些不确定性，我们常把 Hubble 常数参数化为
\begin{equation*}
H_0 = 100\,h\ \mathrm{km/sec/Mpc},
\tag{8.70}
\end{equation*}
于是 $h\approx0.7$。典型的宇宙学尺度由\textbf{Hubble 长度}（Hubble length）设定
\begin{align*}
d_H &= H_0^{-1}c \\
&= 9.25\times10^{27}\,h^{-1}\ \mathrm{cm} \\
&= 3.00\times10^{3}\,h^{-1}\ \mathrm{Mpc},
\tag{8.71}
\end{align*}
以及\textbf{Hubble 时间}（Hubble time）
\begin{align*}
t_H &= H_0^{-1} \\
&= 3.09\times10^{17}\,h^{-1}\ \mathrm{sec} \\
&= 9.78\times10^{9}\,h^{-1}\ \mathrm{yr}.
\tag{8.72}
\end{align*}
当然，由于我们通常取 $c=1$，你会看到 $H_0^{-1}$ 既被称作 Hubble 长度，又被称作 Hubble 时间。此外还有\textbf{减速参数}（deceleration parameter），
\begin{equation*}
q = -\frac{a\ddot{a}}{\dot{a}^{2}},
\tag{8.73}
\end{equation*}
它量度膨胀速率的变化速率。

另一个有用的量是\textbf{密度参数}（density parameter），
\begin{equation*}
\boxed{\,\Omega = \frac{8\pi G}{3H^{2}}\rho = \frac{\rho}{\rho_{\mathrm{crit}}},\,}
\tag{8.74}
\end{equation*}
其中\textbf{临界密度}（critical density）定义为
\begin{equation*}
\rho_{\mathrm{crit}} = \frac{3H^{2}}{8\pi G}.
\tag{8.75}
\end{equation*}
这个量一般会随时间变化，被称为\emph{临界}密度，因为弗里德曼方程 (8.67) 可以写成
\begin{equation*}
\Omega - 1 = \frac{\kappa}{H^{2}a^{2}}.
\tag{8.76}
\end{equation*}
因此，$\kappa$ 的符号由 $\Omega$ 大于、等于还是小于 1 决定。我们有
\begin{align*}
\rho < \rho_{\mathrm{crit}} \quad&\leftrightarrow\quad \Omega < 1 \quad\leftrightarrow\quad \kappa < 0 \quad\leftrightarrow\quad \text{开放} \\
\rho = \rho_{\mathrm{crit}} \quad&\leftrightarrow\quad \Omega = 1 \quad\leftrightarrow\quad \kappa = 0 \quad\leftrightarrow\quad \text{平直} \\
\rho > \rho_{\mathrm{crit}} \quad&\leftrightarrow\quad \Omega > 1 \quad\leftrightarrow\quad \kappa > 0 \quad\leftrightarrow\quad \text{闭合}.
\end{align*}
于是，密度参数告诉我们，三种 Robertson--Walker 几何中哪一种描述我们的宇宙。从观测上确定它是至关重要的；近来对宇宙微波背景各向异性的测量使我们相信，$\Omega$ 非常接近于 1。

% ---- 书页 338 / p353 ----

\section{尺度因子的演化}

给定不同组分 $i$ 的能量密度 $\rho_i$ 的数值，连同它们的物态方程 $p_i=p_i(\rho_i)$，以及空间曲率 $\kappa$ 的大小，就可以求解弗里德曼方程 (8.67)，得到尺度因子 $a(t)$ 演化的完整历史。一般来说，我们只需对弗里德曼方程（它只是一阶微分方程）作数值积分，不过，对不同宇宙学参数所对应的解的类型获得一些直观感受是有用的。

为了简化任务，让我们设想能量密度的各个不同组分都按幂律演化，
\begin{equation*}
\rho_i = \rho_{i0}a^{-n_i}.
\tag{8.77}
\end{equation*}
与式 (8.55) 相比较，这等价于假定每个物态方程参数 $w_i=p_i/\rho_i$ 都是一个常数，其值为
\begin{equation*}
w_i = \frac{1}{3}n_i - 1.
\tag{8.78}
\end{equation*}
把空间曲率的贡献当作一种虚构的能量密度（fictitious energy density）来处理，我们还可以进一步简化表达式：
\begin{equation*}
\rho_{\mathrm{c}} \equiv -\frac{3\kappa}{8\pi G a^{2}},
\tag{8.79}
\end{equation*}
以及相应的密度参数
\begin{equation*}
\Omega_{\mathrm{c}} = -\frac{\kappa}{H^{2}a^{2}}.
\tag{8.80}
\end{equation*}
当然，它\emph{不是}能量密度，所以别忘了这只是一种记号上的障眼法。我们最喜欢的几种源的行为总结在下表中。
\begin{equation*}
\begin{array}{l|c|c}
 & w_i & n_i \\ \cline{2-3}
\text{物质} & 0 & 3 \\
\text{辐射} & \frac{1}{3} & 4 \\
\text{曲率} & -\frac{1}{3} & 2 \\
\text{真空} & -1 & 0
\end{array}
\tag{8.81}
\end{equation*}
在这些变量下，弗里德曼方程 (8.67) 可以写成
\begin{equation*}
H^{2} = \frac{8\pi G}{3}\sum_{i(\mathrm{c})}\rho_i,
\tag{8.82}
\end{equation*}
其中记号 $\sum_{i(\mathrm{c})}$ 表示我们不仅对所有真实的能量密度组分 $\rho_i$ 求和，也对空间曲率的贡献

% ---- 书页 339 / p354 ----
$\rho_{\mathrm{c}}$ 求和。注意，如果把两边除以 $H^{2}$，就得到
\begin{equation*}
1 = \sum_{i(\mathrm{c})}\Omega_i.
\tag{8.83}
\end{equation*}
右边的和\emph{不是}总密度参数 $\Omega$，后者只包含来自真实能量密度（不含曲率）的贡献；因此我们有
\begin{equation*}
\Omega_{\mathrm{c}} = 1 - \Omega.
\tag{8.84}
\end{equation*}

让我们先问一问：如果所有的 $\rho_i$（包括 $\rho_{\mathrm{c}}$）都非负，会发生什么。由于 $H^{2}$ 正比于 $\sum_{i(\mathrm{c})}\rho_i$，只要 $\sum_{i(\mathrm{c})}\rho_i\neq0$，宇宙就永远不会经历从膨胀到收缩的转变。我们还可以对 Hubble 参量求时间导数，
\begin{equation*}
\dot{H} = \frac{\ddot{a}}{a} - \left(\frac{\dot{a}}{a}\right)^{2},
\tag{8.85}
\end{equation*}
并代入两条弗里德曼方程 (8.67) 和 (8.68)，得到
\begin{equation*}
\dot{H} = -4\pi G\sum_{i(\mathrm{c})}(1+w_i)\rho_i.
\tag{8.86}
\end{equation*}
由于我们设想 $|w_i|\le1$，当所有 $\rho_i$ 都非负时，我们将总有 $\dot{H}\le0$。换言之，宇宙持续膨胀，但膨胀速率不断减小（这就引出一个绝佳的问题：是什么让宇宙一开始就如此之大？）。

由式 (8.85) 我们看到，$\ddot{a}$ 可以在 $\dot{H}$ 为负的同时为正——即使以 Hubble 参量量度的膨胀速率在减小，尺度因子也可以在“加速”（例如当 $a\propto t^{2}$ 时）。这是非欧几何不可避免的一个微妙之处。Hubble 参量与尺度因子的导数回答的是两个不同的问题。如果把两个检验粒子放在某个固定的初始距离上，问此后很短的时间内它们分开了多少，答案由 Hubble 参量给出；反之，如果选定某个固定的源，问它随时间看上去如何远离我们，答案则由尺度因子的变化给出。因此，“加速”（或“减速”）有两种截然不同但同样正当的含义。实际上，“加速”通常指 $\ddot{a}>0$ 的情形，即使 $\dot{H}<0$。这番讨论并非纯粹学究式的；正如我们在下文将要看到的，我们当前真实的宇宙似乎正是这种类型。

每个 $\rho_i$ 都非负绝非必然。物质和辐射来自动力学的粒子与场，因此我们预期它们的能量密度永远不会为负；如果它们可以为负，空的空间就可能衰变成一团正、负能量的场。但真空和曲率就是另一回事了。真空能是非动力学的，所以负值不会诱发任何不稳定性；而曲率只是空间几何的一种属性，正负皆可。于是，如果我们有
% ---- 书页 340 / p355 ----
负的真空能或正的空间曲率（记住 $\rho_{\mathrm{c}}\propto-\kappa$），Hubble 参量就可以变为零甚至变号。de Sitter 度规 (8.6) 就提供了一个例子：它既有正的真空能，又有正的空间曲率；它描述了一个先坍缩、到达转折点、然后开始膨胀的宇宙。

真实世界是个不整洁的地方，由许许多多不同种类的能量密度组成。不过，由于不同的源以不同的速率演化，在很长一段时间内能量密度会明显地由某一种源主导。因此，考察弗里德曼方程在只有一种能量密度 $\rho\propto a^{-n}$ 时的解是非常有用的。因为我们把空间曲率也当作一种有效的能量源包括进来，这意味着我们要么考虑由单一源主导的平直宇宙，要么考虑具有空间曲率的完全空无一物的宇宙。弗里德曼方程于是意味着
\begin{equation*}
\dot{a}\propto a^{1-n/2}.
\tag{8.87}
\end{equation*}
这可以立即积分得到
\begin{equation*}
\boxed{\,a\propto t^{2/n}\qquad(\text{当 }\rho\propto a^{-n}\text{ 时}).\,}
\tag{8.88}
\end{equation*}
例如考虑一个由物质主导的平直宇宙，$\Omega=\Omega_{\mathrm{M}}=1$；它被称为 Einstein--de Sitter 模型，并且在很长一段时间内曾是描述真实世界的最受青睐的模型（至少在理论家中间是如此）。在 Einstein--de Sitter 宇宙中，尺度因子按 $a\propto t^{2/3}$ 演化；而平直的辐射主导宇宙则按 $a\propto t^{1/2}$ 演化。任何 $n>2$ 的这类宇宙的共形图在附录 H 中导出。尽管我们相信真实宇宙中物质、辐射和真空能的数量都不为零，这些解仍然非常有用；正如我们后面将要讨论的，宇宙在早期是辐射主导的，而在宇宙从 $a\sim1/3000$ 膨胀到 $a\sim1/2$ 的过程中则是物质主导的。

这些解都有一个 $a=0$ 处的奇点，称为\textbf{大爆炸}（Big Bang）。它代表宇宙从奇性状态的创生，而不是物质爆炸进入一个先前存在的时空。人们或许会希望，我们的 FRW 宇宙的完美对称性正是这个奇点的起因，但事实并非如此；宇宙学奇点定理（cosmological singularity theorems）表明，任何满足 $\rho>0$ 且 $p\ge0$ 的宇宙都必定始于一个奇点。当然，当 $a\to0$ 时能量密度变得任意高，我们不指望经典广义相对论在这个区域还能准确地描述自然；想必量子引力将变得重要，尽管目前尚不清楚它如何起作用。

看式 (8.88) 可见，真空能主导（$n=0$）的宇宙显然是一个特例。此时尺度因子按指数而非幂律膨胀；整个度规为
\begin{equation*}
\mathrm{d}s^{2} = -\mathrm{d}t^{2} + e^{Ht}[\mathrm{d}x^{2} + \mathrm{d}y^{2} + \mathrm{d}z^{2}],
\tag{8.89}
\end{equation*}

% ---- 书页 341 / p356 ----
其中 Hubble 参量 $H$ 是常数。当然，在 8.1 节中我们已经描述过一个具有正宇宙学常数的宇宙学时空：de Sitter 空间，其特征是 $\kappa>0$ 且 $a\propto\cosh(t/\alpha)$。那个解与此处 $\kappa=0$、$a\propto\exp(Ht)$ 的解之间是什么关系？它们是同一个时空，只是用不同的坐标来表示。验证这一点的一种办法，是计算 (8.89) 的 Riemann 张量，并检验它具有最大对称时空的特征形式，即式 (8.1)。由于正曲率的最大对称时空在局部是唯一的，度规 (8.6) 与 (8.89) 必定描述同一个流形，或其一部分。实际上，(8.89) 的坐标只覆盖了 de Sitter 空间的一部分；它们在过去是不完备的。在习题中你会被要求证明，这些坐标中的共动测地线会在有限的仿射参量内到达 $t=-\infty$；它们撞上了坐标的边缘。在图 8.1 的共形图中，这些坐标覆盖了正方形右上方的三角形部分。关于 de Sitter 空间和反 de Sitter 空间上各种坐标系的更完备的描述，参见 Hawking and Ellis (1973)。

另一个有趣的特例是 $\rho=0$ 但具有空间曲率的完全空无一物的宇宙。弗里德曼方程变为
\begin{equation*}
H^{2} = -\frac{\kappa}{a^{2}},
\tag{8.90}
\end{equation*}
所以曲率 $\kappa$ 必须为负。把曲率看作一种虚构的能量密度 $\rho_{\mathrm{c}}\propto a^{-2}$，由式 (8.88) 可知，这样的宇宙将线性膨胀，$a\propto t$。这个时空被称为\textbf{Milne 宇宙}（Milne universe）。然而，正如 de Sitter 空间的情形一样，我们知道另一个 $\rho=0$ 的宇宙学时空——在这里就是平直的 Minkowski 空间。再一次，Milne 时空只是 Minkowski 空间在某个不完备坐标系中的一个片。它可以被看作 Minkowski 空间中某个固定点的未来光锥的内部，用负曲率的双曲面来分层。要检验这一点，只需计算 Riemann 张量的所有分量——结果它们全部为零；任何 Riemann 曲率为零的时空在局部都是 Minkowski 空间。

与这些理想化的解相反，一个现实的宇宙学将包含好几种形式的能量-动量。在当前的宇宙中，我们有把握认为辐射密度显著低于物质密度，但真空与物质在动力学上都很重要。因此，方便的做法是用 $\Omega_{\mathrm{M}}$ 和 $\Omega_{\Lambda}$ 来参数化像我们这样的宇宙，曲率则由 $\Omega_{\mathrm{c}}=1-\Omega_{\mathrm{M}}-\Omega_{\Lambda}$ 固定。这类宇宙某些具体例子的膨胀历史示于图 8.3。随着这些宇宙的膨胀，物质、曲率和真空的相对影响会发生变化，因为相应的密度以不同的速率演化：
\begin{equation*}
\Omega_{\Lambda}\propto\Omega_{\mathrm{c}}a^{2}\propto\Omega_{\mathrm{M}}a^{3}.
\tag{8.91}
\end{equation*}
当 $a\to0$ 的过去，曲率和真空将可以忽略，宇宙的行为如同 Einstein--de Sitter。当 $a\to\infty$ 的未来，曲率和物质将可以忽略，宇宙将渐近于 de Sitter 空间；除非尺度因子
% ---- 书页 342 / p357 ----
永远达不到无限大，因为宇宙会在某个有限的时刻开始再度坍缩。

\begin{figure}[H]
\centering
\includegraphics[width=0.72\textwidth]{fig_8.3.pdf}
\caption*{图 8.3\quad 不同 $\Omega_{\mathrm{M}}$ 与 $\Omega_{\Lambda}$ 取值下的膨胀历史。自上而下，各条曲线分别描述 $(\Omega_{\mathrm{M}},\Omega_{\Lambda})=(0.3,0.7)$、$(0.3,0.0)$、$(1.0,0.0)$ 和 $(4.0,0.0)$。}
\end{figure}

如果真空能为负，就\emph{总会}发生再度坍缩；随着宇宙膨胀，真空能最终将占据主导，而 $\Omega_{\Lambda}<0$ 的效果是导致减速和再度坍缩（正如 $\Omega_{\Lambda}>0$ 的效果是把宇宙推开）。当 $\Omega_{\Lambda}\ge0$ 时，如果 $\Omega_{\mathrm{M}}$ 大到足以在 $\Omega_{\Lambda}$ 来得及接管之前使宇宙的膨胀停下来，再度坍缩也是可能的。这些可能性表现为图 8.4 中 $\Omega_{\mathrm{M}}/\Omega_{\Lambda}$ 参数空间的不同区域。对角线代表 $\Omega_{\mathrm{total}}=1$，意味着 $\kappa=0$。

要确定永恒膨胀与最终再度坍缩之间的分界线，请注意：坍缩要求 Hubble 参量在从正变负的过程中经过零。发生这种转折处的尺度因子 $a_{*}$，可以通过在弗里德曼方程中令 $H=0$ 而求得，
\begin{equation*}
H^{2} = 0 = \frac{8\pi G}{3}\left(\rho_{\mathrm{M}0}a_{*}^{-3} + \rho_{\Lambda0} + \rho_{\mathrm{c}0}a_{*}^{-2}\right).
\tag{8.92}
\end{equation*}
把它除以 $H_0^{2}$，利用 $\Omega_{\mathrm{c}0}=1-\Omega_{\mathrm{M}0}-\Omega_{\Lambda0}$，再稍加整理，便得到
\begin{equation*}
\Omega_{\Lambda0}a_{*}^{3} + (1-\Omega_{\mathrm{M}0}-\Omega_{\Lambda})a_{*} + \Omega_{\mathrm{M}0} = 0.
\tag{8.93}
\end{equation*}
（译注：按式 (8.92) 的推导，第二项应为 $\Omega_{\Lambda0}$；原书排印漏写下标 0，原文如此。）

这是关于 $a_{*}$（转折点处的尺度因子）的三次方程。当然，我们对 $a_{*}$ 本身并不很在乎；我们在乎的是，给定 $\Omega_{\mathrm{M}0}$，使 (8.93) 存在实解的 $\Omega_{\Lambda0}$ 值。求解这个三次方程并做一些数学运算，我们发现，宇宙将永远膨胀所对应的 $\Omega_{\Lambda0}$ 值由下式给出

% ---- 书页 343 / p358 ----
\begin{figure}[H]
\centering
\includegraphics[width=0.72\textwidth]{fig_8.4.pdf}
\caption*{图 8.4\quad 由物质和真空能主导的宇宙的性质，表示为密度参数 $\Omega_{\mathrm{M}}$ 与 $\Omega_{\Lambda}$ 的函数。左上角的圆形区域大致代表实验数据（截至 2003 年）所偏好的数值。}
\end{figure}

\begin{equation*}
\Omega_{\Lambda0} \ge
\left\{
\begin{array}{ll}
0 & 0\le\Omega_{\mathrm{M}0}\le1 \\
4\Omega_{\mathrm{M}0}\cos^{3}\left[\dfrac{1}{3}\cos^{-1}\left(\dfrac{1-\Omega_{\mathrm{M}0}}{\Omega_{\mathrm{M}0}}\right) + \dfrac{4\pi}{3}\right] & \Omega_{\mathrm{M}0}>1.
\end{array}
\right.
\tag{8.94}
\end{equation*}
注意，当 $\Omega_{\Lambda0}=0$ 时，开放和平直宇宙（$\Omega_0=\Omega_{\mathrm{M}0}\le1$）将永远膨胀，而闭合宇宙（$\Omega_0=\Omega_{\mathrm{M}0}>1$）将再度坍缩。对宇宙学常数由来已久的轻视，造成了一种民间信念，认为这是一种必然的对应关系；然而，一旦承认真空能的可能性，空间几何与最终命运之间的任何组合就都是可能的了。

在图 8.4 的左上角，我们标出了目前受偏好的宇宙学参数取值：$\Omega_{\mathrm{M}0}\sim0.3$、$\Omega_{\Lambda0}\sim0.7$，我们将在 8.7 节中讨论。这已深入永恒膨胀的区域；如果真空能保持真正的恒定（它未必如此），我们的宇宙就注定要永远膨胀下去。

在本节的结尾，我们指出为弗里德曼方程寻找静态解的困难。要成为静态，我们必须不仅有 $\dot{a}=0$，还要 $\ddot{a}=0$。由式 (8.68)，这只有在压强为
\begin{equation*}
p = -\frac{1}{3}\rho,
\tag{8.95}
\end{equation*}
时才可能；而由式 (8.67)，还必须存在非零的空间曲率
\begin{equation*}
\frac{\kappa}{a^{2}} = \frac{8\pi G}{3}\rho.
\tag{8.96}
\end{equation*}

% ---- 书页 344 / p359 ----
由于能量密度和压强必须异号，如果我们只诉诸物质或辐射，这些条件就无法满足。当 Einstein 最初在广义相对论中寻找宇宙学解时，天文学家尚未发现宇宙正在膨胀，因此静态解的缺失曾被认为是个难题。这为 Einstein 引入宇宙学常数提供了动机；静态条件可以由物质和真空能的组合来满足，其中
\begin{equation*}
\rho_{\Lambda} = \frac{1}{2}\rho_{\mathrm{M}},
\tag{8.97}
\end{equation*}
再加上适当的正空间曲率。这些参数描述的是\textbf{Einstein 静态宇宙}（Einstein static universe）。如今我们知道宇宙在膨胀，因此这个解没有多少经验上的意义；但它对理论家却极为有用，为共形图的构造提供了基础。

\section{红移与距离}

显然，我们希望通过观测测定若干个量，以判断 FRW 模型中的哪一个对应我们的宇宙。显然我们想测定 $H_0$，因为它与宇宙的年龄有关；我们还想知道 $\Omega$，它通过式 (8.76) 决定 $\kappa$。为了理解这些量究竟可能如何被测量，让我们考虑 FRW 宇宙中的测地线运动。FRW 宇宙中存在若干类空 Killing 矢量，却没有能给我们守恒能量概念的类时 Killing 矢量。不过，确实存在一个 Killing 张量。如果 $U^{\mu}=(1,0,0,0)$ 是共动观者的四速度，那么张量
\begin{equation*}
K_{\mu\nu} = a^{2}(g_{\mu\nu} + U_{\mu}U_{\nu})
\tag{8.98}
\end{equation*}
满足 $\nabla_{(\sigma}K_{\mu\nu)}=0$（你可以自行验证），因而是一个 Killing 张量。这意味着，如果粒子具有四速度 $V^{\mu}=\mathrm{d}x^{\mu}/\mathrm{d}\lambda$，那么
\begin{equation*}
K^{2} = K_{\mu\nu}V^{\mu}V^{\nu} = a^{2}[V_{\mu}V^{\mu} + (U_{\mu}V^{\mu})^{2}]
\tag{8.99}
\end{equation*}
沿测地线将是常数。让我们先就有质量粒子的情形来思考这个问题。此时我们将有 $V_{\mu}V^{\mu}=-1$，于是
\begin{equation*}
(V^{0})^{2} = 1 + |\vec{V}|^{2},
\tag{8.100}
\end{equation*}
其中 $|\vec{V}|^{2}=g_{ij}V^{i}V^{j}$。我们还有 $U_{\mu}V^{\mu}=-V^{0}$，所以式 (8.99) 给出
\begin{equation*}
|\vec{V}| = \frac{K}{a}.
\tag{8.101}
\end{equation*}
因此，随着宇宙膨胀，粒子相对于共动坐标会“慢下来”。事实上这是一种真实的变慢，其意义是：初始相对速度很高的一团粒子气体将随宇宙膨胀而冷却下来。

% ---- 书页 345 / p360 ----
类似的事情也发生在类光测地线上。此时 $V_{\mu}V^{\mu}=0$，式 (8.99) 给出
\begin{equation*}
U_{\mu}V^{\mu} = \frac{K}{a}.
\tag{8.102}
\end{equation*}
但共动观者测得的光子频率是 $\omega=-U_{\mu}V^{\mu}$。因此，以频率 $\omega_{\mathrm{em}}$ 发出的光子，随着宇宙的膨胀将被观测到较低的频率 $\omega_{\mathrm{obs}}$：
\begin{equation*}
\frac{\omega_{\mathrm{obs}}}{\omega_{\mathrm{em}}} = \frac{a_{\mathrm{em}}}{a_{\mathrm{obs}}}.
\tag{8.103}
\end{equation*}
宇宙学家喜欢用这两个事件之间的\textbf{红移}（redshift）$z$ 来谈论它，红移由波长的分数变化定义：
\begin{equation*}
z_{\mathrm{em}} = \frac{\lambda_{\mathrm{obs}} - \lambda_{\mathrm{em}}}{\lambda_{\mathrm{em}}}.
\tag{8.104}
\end{equation*}
如果观测发生在今天（$a_{\mathrm{obs}}=a_0=1$），这意味着
\begin{equation*}
\boxed{\,a_{\mathrm{em}} = \frac{1}{1+z_{\mathrm{em}}}.\,}
\tag{8.105}
\end{equation*}
所以，一个天体的红移告诉我们光子发射时的尺度因子。

注意，这种红移与传统的多普勒效应并不是一回事；导致红移的是空间的膨胀，而不是观者与发射者的相对速度。不过，如果我们在与 Hubble 半径 $H_0^{-1}$ 以及空间曲率半径 $\kappa^{-1/2}$ 相比很小的距离上观测星系，宇宙的膨胀看上去就很像一群星系在彼此远离，红移看上去也很像多普勒效应。因此，天文学家常常用“速度”$v=cz$（$c$ 是光速）来考虑红移。尽管我们知道，对弯曲时空中不同点上的两个物体并不能真正谈论相对速度，但这个虚构的说法在足够短的距离上工作得很好。在这个近似之内，从我们到某个星系的“距离”$d$ 可以取为\textbf{瞬时物理距离}（instantaneous physical distance）$d_P$（即在诸如厘米这样的物理单位下，沿我们当前的空间超曲面量出的、我们与该星系位置之间的距离）。让我们把 RW 度规写成
\begin{equation*}
\mathrm{d}s^{2} = -\mathrm{d}t^{2} + a^{2}(t)R_0^{2}\left[\mathrm{d}\chi^{2} + S_k^{2}(\chi)\,\mathrm{d}\Omega^{2}\right],
\tag{8.106}
\end{equation*}
其中 $S_k(\chi)$ 由式 (8.33) 定义，而 $k\in\{+1,0,-1\}$。在这种形式下，在时刻 $t$ 测得的、我们（$\chi=0$）与共动径向坐标为 $\chi$ 的星系之间的瞬时物理距离就是
\begin{equation*}
d_P(t) = a(t)R_0\chi,
\tag{8.107}
\end{equation*}

% ---- 书页 346 / p361 ----
其中 $\chi$ 保持不变，因为我们假定我们和被观测的星系都完美地共动。（它们可能并不如此；在那种情形下，把所谓“本动速度”（peculiar velocities）引起的修正包括进来是轻而易举的。）当然，“距离”之所以加引号，是因为一旦离开这个近似，就有好几种彼此不等价而有用的距离概念；但当 $d_P$ 很小时它们全都一致。这时，观测到的速度（由红移推断）就简单地是
\begin{equation*}
v = \dot{d}_P = \dot{a}R_0\chi = \frac{\dot{a}}{a}d_P.
\tag{8.108}
\end{equation*}
在今天取值，它就成为
\begin{equation*}
\boxed{\,v = H_0 d_P,\,}
\tag{8.109}
\end{equation*}
这就是著名的\textbf{Hubble 定律}（Hubble law）：对于不太遥远的星系，观测到的退行速度与距离成正比。

如果红移不是非常小，我们就得更仔细地思考宇宙学中“距离”指的是什么。瞬时物理距离是一个方便的构造，但它本身不可观测，因为观测涉及的总是我们过去光锥上的事件，而不是我们当前的空间超曲面。在欧几里得空间中，推断一个物体的距离有许多不同的方法；例如，我们可以把它的视亮度与它的内禀光度相比较，或者把它的视角速度与它的内禀横向速度相比较，或者把它的视角大小与它的物理尺度相比较。对每一种情形，我们都可以定义一种距离，即假如空间是欧几里得的、宇宙不在膨胀时我们\emph{将会}推断出的距离。

让我们从\textbf{光度距离}（luminosity distance）$d_L$ 开始，它的定义是满足
\begin{equation*}
d_L^{2} = \frac{L}{4\pi F},
\tag{8.110}
\end{equation*}
其中 $L$ 是源的绝对光度，$F$ 是观者测得的流量（flux），即单位时间落到某个探测器单位面积上的能量。这个定义来自如下事实：在平直空间中，对位于距离 $d$ 处的源，流量与光度之比正好是以源为中心的球面面积的倒数，$F/L = 1/A(d) = 1/4\pi d^{2}$。然而在 FRW 宇宙中，流量将被稀释。光子数守恒告诉我们，源发出的所有光子最终都将穿过与发射者相距共动距离 $\chi$ 的球面。但流量还会被两个附加效应稀释：单个光子红移一个因子 $(1+z)$；而且光子撞击球面的频率变低，因为相隔 $\delta t$ 发出的两个光子，被观测到的相隔时间是 $(1+z)\delta t$。于是我们将有
\begin{equation*}
\frac{F}{L} = \frac{1}{(1+z)^{2}A}.
\tag{8.111}
\end{equation*}

% ---- 书页 347 / p362 ----
以共动距离 $\chi$ 为中心的球面的面积 $A$，可以从式 (8.106) 中 $\mathrm{d}\Omega^{2}$ 的系数导出，得到
\begin{equation*}
A = 4\pi R_0^{2}S_k^{2}(\chi),
\tag{8.112}
\end{equation*}
其中我们已取 $a(t)=1$，因为我们是在今天观测这些光子。把所有这些合在一起，就得到
\begin{equation*}
d_L = (1+z)R_0S_k(\chi).
\tag{8.113}
\end{equation*}
光度距离 $d_L$ 是我们有希望测量的量，因为有一些天体物理源的绝对光度是已知的。但 $\chi$ 是不可观测的，所以我们必须把它从方程中消去。在类光测地线（为方便起见取为径向）上，我们有
\begin{equation*}
0 = \mathrm{d}s^{2} = -\mathrm{d}t^{2} + a^{2}R_0^{2}\mathrm{d}\chi^{2},
\tag{8.114}
\end{equation*}
即
\begin{equation*}
\chi = R_0^{-1}\int\frac{\mathrm{d}t}{a} = R_0^{-1}\int\frac{\mathrm{d}a}{a^{2}H(a)},
\tag{8.115}
\end{equation*}
其中我们用了 $H=\dot{a}/a$。惯常的做法是用 $a=1/(1+z)$ 把尺度因子换成红移，于是我们有
\begin{equation*}
\chi(z) = R_0^{-1}\int_{0}^{z}\frac{\mathrm{d}z'}{H(z')}.
\tag{8.116}
\end{equation*}

为了计算这个积分中的 Hubble 参量，我们使用弗里德曼方程 (8.67)，像上一节中那样把它写成
\begin{equation*}
H^{2} = \frac{8\pi G}{3}\sum_{i(\mathrm{c})}\rho_i.
\tag{8.117}
\end{equation*}
为了简化问题，我们可以再次假设每个密度组分都按幂律演化，
\begin{equation*}
\rho_i(z) = \rho_{i0}a^{-n_i} = \rho_{i0}(1+z)^{n_i},
\tag{8.118}
\end{equation*}
于是我们可以写出
\begin{equation*}
H(z) = H_0E(z),
\tag{8.119}
\end{equation*}
其中
\begin{equation*}
E(z) = \left[\sum_{i(\mathrm{c})}\Omega_{i0}(1+z)^{n_i}\right]^{1/2},
\tag{8.120}
\end{equation*}

% ---- 书页 348 / p363 ----
其中密度参数 $\Omega_i$ 由式 (8.74) 定义。下面涉及 $E(z)$ 的那些方程，无论能量源是否按幂律演化都成立；若不按幂律演化，只需采用 $E(z)=H(z)/H_0$［其中 $H(z)$ 由弗里德曼方程确定］而不必用式 (8.120)。

于是，光度距离为
\begin{equation*}
d_L(z) = (1+z)R_0S_k\left[R_0^{-1}H_0^{-1}\int\frac{\mathrm{d}z'}{E(z')}\right].
\tag{8.121}
\end{equation*}
注意，当 $k=0$ 时 $R_0$ 自动消去，这很好，因为在那种情形下它是一个完全任意的参数。即使它并非任意，人们通常也更愿意用 $\Omega_{\mathrm{c}0}=-k/R_0^{2}H_0^{2}$ 来讨论；这个量既可以由对空间曲率的测定直接测量，也可以通过测量密度参数并利用 $\Omega_{\mathrm{c}0}=1-\Omega_0$ 来得到。用这个参数，我们有
\begin{equation*}
R_0 = H_0^{-1}\sqrt{-k\Omega_{\mathrm{c}0}} = \frac{H_0^{-1}}{\sqrt{|\Omega_{\mathrm{c}0}|}}.
\tag{8.122}
\end{equation*}
于是，我们把光度距离用可测量的宇宙学参数写成
\begin{equation*}
\boxed{\,d_L(z) = (1+z)\frac{H_0^{-1}}{\sqrt{|\Omega_{\mathrm{c}0}|}}S_k\left[\sqrt{|\Omega_{\mathrm{c}0}|}\int\frac{\mathrm{d}z'}{E(z')}\right].\,}
\tag{8.123}
\end{equation*}
这个方程虽然显得笨拙，却在宇宙学中处于核心地位。给定可观测量 $H_0$ 和 $\Omega_{i0}$，我们可以直接算出到位于任一红移 $z$ 处天体的光度距离；同样地，我们也可以对一系列红移范围内的天体测量 $d_L(z)$，并从这些信息中提取 $H_0$ 和/或各个 $\Omega_{i0}$。

除光度距离之外，还有两个相关的距离量度。正如光度距离是假如我们身处平直空间时，从源的内禀光度和观测光度推断出的距离，\textbf{固有运动距离}（proper motion distance）$d_M$ 是从源的内禀运动和观测运动推断出的距离。它的定义为
\begin{equation*}
d_M = \frac{u}{\dot{\theta}},
\tag{8.124}
\end{equation*}
其中 $u$ 是固有横向速度（比如你会以 k/s 为单位来测量的量），而 $\dot{\theta}$ 是观测到的角速度。另一方面，\textbf{角直径距离}（angular diameter distance）是从源的内禀大小和观测大小推断出的距离；它的定义为
\begin{equation*}
d_A = \frac{R}{\theta},
\tag{8.125}
\end{equation*}

% ==== tail：8.5 节（红移与距离）收尾，承接书页 348 末的 (8.125)，
% ==== 至书页 349 中 8.6 节标题之前；置于本文件第一个 \section 之前 ====

% ---- 书页 349 / p364 ----
%JOIN%其中 $R$ 是天体的固有大小，$\theta$ 是观测到的角直径。在这两种情形下，我们都能导出与式 (8.123) 类似的公式；所幸，对宇宙学参数的那种难以驾驭的依赖关系为所有距离量度所共有，最后剩下的只是对红移的简单依赖：
\begin{equation*}
d_L = (1+z)d_M = (1+z)^{2}d_A,
\tag{8.126}
\end{equation*}
鼓励各位动手验证一下。这样，只要测得了其中一种距离，我们就能方便地换算成其他任何一种；或者，我们也可以独立地测量不同的距离，并利用式 (8.126) 来检验 RW 框架的自洽性。

趁我们还在思量距离，不妨也考虑一下：从红移为 $z$ 的天体发出的光，从发射时刻到现在流逝了多长时间。如果宇宙今天的年龄是 $t_0$，而光子被发射时宇宙的年龄是 $t_{*}$，那么\textbf{回溯时间}（lookback time）为
\begin{align*}
t_0 - t_{*} &= \int_{t_{*}}^{t_0}\mathrm{d}t\\
&= \int_{a_{*}}^{1}\frac{\mathrm{d}a}{aH(a)}\\
&= H_0^{-1}\int_0^{z_{*}}\frac{\mathrm{d}z'}{(1+z')E(z')}.
\tag{8.127}
\end{align*}
例如，考虑一个平直（$k=0$）、物质主导（$\rho=\rho_{\mathrm{M}}=\rho_{\mathrm{M}0}a^{-3}$）的宇宙。那么
\begin{equation*}
E(z) = (1+z)^{3/2},
\tag{8.128}
\end{equation*}
于是
\begin{align*}
t_0 - t_{*} &= H_0^{-1}\int_0^{z_{*}}\frac{\mathrm{d}z'}{(1+z')^{5/2}}\\
&= \frac{2}{3}H_0^{-1}\left[1-(1+z_{*})^{-3/2}\right].
\tag{8.129}
\end{align*}
令 $t_{*}\to 0$（$z_{*}\to\infty$），便得到物质主导宇宙的总年龄：
\begin{equation*}
t_0(\mathrm{MD}) = \frac{2}{3}H_0^{-1}.
\tag{8.130}
\end{equation*}
对于并非完全物质主导的宇宙，$\frac{2}{3}$ 这个因子不会完全正确，但对于合理的宇宙学参数取值，我们通常得到 $t_0\sim H_0^{-1}$。

\section{引力透镜}

在第 7 章中我们介绍了引力透镜（gravitational lensing）的概念：光在 Newton 引力场中的偏折与时间延迟。除了在太阳系内为广义相对论提供检验之外，透镜效应还出现在众多天体物理情境之中，并已成为现代宇宙学不可或缺的组成部分。\footnote{关于引力透镜的一篇优秀综述——我们的讨论即从中借鉴——参见 R.~Narayan and M.~Bartelmann, ``Lectures on Gravitational Lensing,'' 13th Jerusalem Winter School in Theoretical Physics, \texttt{http://arxiv.org/astro-ph/9606001}。}

% ---- 书页 350 / p365 ----
\begin{figure}[H]
\centering
\includegraphics[width=0.72\textwidth]{fig_8.5.pdf}
\caption*{图 8.5\quad 引力透镜的几何，概括于透镜方程 (8.132) 之中。透镜的作用，是把平直 Minkowski 背景中本应观测到的角 $\beta$ 扭变为角 $\theta$。}
\end{figure}

宇宙学透镜效应与我们先前讨论的情形有两个重要区别：Robertson--Walker 度规取代了 Minkowski 背景，而且透镜本身往往也比简单的点质量复杂。图 8.5 描绘了一种典型的透镜几何。在这一节的讨论中，我们将始终假设透镜是“薄”的——即其空间延展尺度远小于源、透镜与观者之间的距离。在这种情形下，我们可以合理地谈论一个唯一的到透镜的距离 $d_L$，以及透镜到源的距离 $d_{LS}$。

对于天空上一幅（可能很复杂的）像，我们用像的不同成分之间的一组夹角来描述它。这些夹角可以看作天空上的二维矢量。透镜的作用，是把没有任何偏折时本应观测到的那些角——例如源与透镜之间的夹角 $\vec{\beta}$——扭变为一幅由一组角 $\vec{\theta}$ 刻画的新像。我们全程假设这些角都很小。这一映射由\textbf{约化透镜角}（reduced lensing angle）$\vec{\alpha}=\vec{\theta}-\vec{\beta}$ 描述。按照图 8.5 所示的几何，它与真实偏转角 $\widehat{\vec{\alpha}}$ 之间的关系为
\begin{equation*}
\vec{\alpha} = \frac{d_{LS}}{d_S}\widehat{\vec{\alpha}}.
\tag{8.131}
\end{equation*}

% ---- 书页 351 / p366 ----
于是我们得到\textbf{透镜方程}（lens equation）
\begin{equation*}
\boxed{\,\vec{\beta} = \vec{\theta} - \frac{d_{LS}}{d_S}\widehat{\vec{\alpha}}.\,}
\tag{8.132}
\end{equation*}
透镜方程所描述的，无非是受扰时空中的光线追踪。

当然，我们应当仔细琢磨图中画出的那些“距离” $d_i$。透镜效应发生在膨胀的宇宙中，而宇宙还可能带有空间曲率。不过，只要我们\emph{定义}距离 $d_i$ 使得透镜方程所描述的几何关系得以成立，透镜方程就依然成立。换言之，这些距离是：在静态欧几里得空间背景中，给定夹角与横向物理尺度时我们本会推断出的距离。而这恰恰就是角直径距离 (8.125) 的定义。因此，本节中的所有距离都取为角直径距离。注意，角直径距离未必能够相加，故 $d_S\neq d_L+d_{LS}$。

举一个简单的例子：点质量透镜。在第 7 章研究 Newton 极限时我们曾经发现，光子穿过引力势 $\Phi$ 时的偏转角由下式给出：
\begin{equation*}
\widehat{\vec{\alpha}} = 2\int \vec{\nabla}_{\perp}\Phi\,\mathrm{d}s,
\tag{8.133}
\end{equation*}
对于碰撞参数为 $b$ 的点质量 $M$，上式化为
\begin{equation*}
\widehat{\vec{\alpha}} = \frac{4GM}{b}.
\tag{8.134}
\end{equation*}
碰撞参数可以写成 $b=d_L\theta$。透镜方程 (8.132) 便成为
\begin{equation*}
\beta = \theta - \frac{d_{LS}}{d_S d_L}\frac{4GM}{\theta}.
\tag{8.135}
\end{equation*}
富有启发性的做法是考虑最简单的情形：源与透镜共线（$\beta=0$）。此时源将被透镜成一个环绕透镜的\textbf{Einstein 环}（Einstein ring），其角距离由\textbf{Einstein 角}（Einstein angle）给出：
\begin{equation*}
\theta_{\mathrm{E}} = \sqrt{\frac{4GMd_{LS}}{d_L d_S}}.
\tag{8.136}
\end{equation*}
即便在更复杂的位形中，Einstein 角也为透镜效应设定了一个特征尺度。我们还可以定义一个与之相联系的距离尺度，即\textbf{Einstein 半径}（Einstein radius）：
\begin{equation*}
R_{\mathrm{E}} = \sqrt{\frac{4GMd_L d_{LS}}{d_S}}.
\tag{8.137}
\end{equation*}

% ---- 书页 352 / p367 ----
换算成厘米或其他物理单位时，别忘了在我们的所有方程中都有 $c=1$。为了对典型天体物理情形下透镜效应的大小有些感觉，我们可以考虑两种常见的场合：我们银河系内近似太阳质量的天体造成的“微透镜”（microlensing），以及星系或星系团造成的宇宙学透镜。前一情形下 Einstein 角将是毫角秒的量级，而后一情形下则是角秒的量级：
\begin{align*}
\theta_{\mathrm{E}} &= 0.9\sqrt{\left(\frac{M}{M_{\odot}}\right)\left(\frac{10\ \mathrm{kpc}}{D}\right)}\ \text{毫角秒}\\
&= 0.9\sqrt{\left(\frac{M}{10^{11}M_{\odot}}\right)\left(\frac{\mathrm{Gpc}}{D}\right)}\ \text{角秒}.
\tag{8.138}
\end{align*}
我们暂时继续讨论点质量透镜。尽管已经观测到若干例壮观的 Einstein 环，但多数时候我们不会幸运到源与透镜恰好完美对齐。这时我们可以求解 (8.135)，得到像角的两个数值：
\begin{equation*}
\theta_{\pm} = \frac{1}{2}\left(\beta \pm \sqrt{\beta^{2}+4\theta_{\mathrm{E}}^{2}}\right).
\tag{8.139}
\end{equation*}
位于 $\theta_{+}$ 处的像总在 Einstein 角之外，而 $\theta_{-}$ 处的像则在其内。实际上这个公式有些误导性：像的数目永远会是奇数；对点质量透镜而言，第三个像将位于与透镜自身相同的位置。

现在来考虑比点质量更一般的透镜。我们知道，偏转角可以借助 Newton 引力势由式 (8.133) 给出。我们可以沿从观者出发指向过去的测地线路径作积分，定义\textbf{透镜势}（lensing potential）：
\begin{equation*}
\psi(\vec{\theta}) = 2\frac{d_{LS}}{d_L d_S}\int \Phi(d_L\vec{\theta},s)\,\mathrm{d}s.
\tag{8.140}
\end{equation*}
利用透镜势，只需取梯度便可直接导出约化透镜角：
\begin{align*}
\vec{\alpha} &= \vec{\nabla}_{\theta}\psi\\
&= 2\frac{d_{LS}}{d_S}\int \vec{\nabla}_{\perp}\Phi\,\mathrm{d}s.
\tag{8.141}
\end{align*}
注意，角梯度 $\vec{\nabla}_{\theta}$ 与 $\vec{\nabla}_{\perp}$（对透镜所在处横向距离的梯度）相差一个因子 $d_L$。薄透镜近似使我们可以把积分收缩为在透镜位置上取值的量。我们还可以对透镜势作用（二维）拉普拉斯算符，得到\textbf{会聚}（convergence）$\kappa$：
\begin{align*}
\kappa(\vec{\theta}) &\equiv \frac{1}{2}\nabla_{\theta}^{2}\psi\\
&= \frac{d_L d_{LS}}{d_S}\int \nabla^{2}\Phi\,\mathrm{d}s.
\tag{8.142}
\end{align*}

% ---- 书页 353 / p368 ----
会聚可以被看作对积分质量密度的一种量度。我们可以把上面的表达式反过来，用会聚同时写出透镜势与约化偏转角：
\begin{equation*}
\psi(\vec{\theta}) = \frac{1}{\pi}\int \kappa(\vec{\theta}\,')\ln|\vec{\theta}-\vec{\theta}\,'|\,\mathrm{d}^{2}\theta'
\tag{8.143}
\end{equation*}
以及
\begin{equation*}
\vec{\alpha}(\vec{\theta}) = \frac{1}{\pi}\int \kappa(\vec{\theta}\,')\frac{\vec{\theta}-\vec{\theta}\,'}{|\vec{\theta}-\vec{\theta}\,'|}\,\mathrm{d}^{2}\theta'.
\tag{8.144}
\end{equation*}
要检验这些方程，请记住这些矢量只定义在两个横向维度上。

会聚描述的是引力透镜对光线的聚焦。这种聚焦使源显得更大（就像放大镜一样）。根据 Liouville 定理——源所发射光子的相空间密度守恒——源的面亮度在透镜作用下保持不变；因此尺度的增大导致亮度的放大。与此同时，光线穿过透镜时的扭转还会引起畸变，导致像的形状发生切变。为了描述这两种现象，我们考虑透镜映射的导数所构成的 $2\times 2$ 矩阵
\begin{equation*}
A_{ij} \equiv \frac{\partial\beta^{i}}{\partial\theta^{j}}.
\tag{8.145}
\end{equation*}
注意，上、下指标之间实际上没有区别，因为它们定义在二维欧几里得平面上。由于 $\vec{\beta}=\vec{\theta}-\vec{\alpha}$，我们有
\begin{align*}
A_{ij} &= \delta_{ij} - \frac{\partial\alpha^{i}}{\partial\theta^{j}}\\
&= \delta_{ij} - \psi_{ij},
\tag{8.146}
\end{align*}
这里我们引入了记号
\begin{equation*}
\psi_{ij} \equiv \frac{\partial^{2}\psi}{\partial\theta^{i}\partial\theta^{j}}.
\tag{8.147}
\end{equation*}
矩阵 $A$ 编码了透镜映射的局域性质。它的逆矩阵被称为\emph{放大张量}（magnification tensor）
\begin{equation*}
\boldsymbol{M} = \frac{\partial\vec{\theta}}{\partial\vec{\beta}} = A^{-1}.
\tag{8.148}
\end{equation*}

% ---- 书页 354 / p369 ----
这个名字从何而来？透镜把由 $\vec{\beta}$ 描述的面元扭变为由 $\vec{\theta}$ 描述的面元，面积的改变由这个映射的雅可比行列式描述，而后者无非是 $\boldsymbol{M}$ 的行列式。这个行列式被定义为\textbf{放大率}（magnification）$\mu$：
\begin{equation*}
\mu = |\boldsymbol{M}| = \frac{1}{|A|}.
\tag{8.149}
\end{equation*}
$\mu$ 的绝对值告诉我们源的亮度实际改变了多少；$\mu$ 可能为负，这意味着像的奇偶性被翻转了。我们说“放大”，是因为只有当透镜和源在天空中彼此靠近时透镜效应才会惹人注意，此时聚焦效应只会带来视亮度的增加；而一个远离源（在天空中位置相距甚远）的透镜所导致的光度极微小的减少，则永远不会被注意到。（如果存在多重像，所有像的亮度之和将超过未畸变源的亮度。）

$A$ 的分量可以分解为会聚与切变的效应。对于会聚，由 $\kappa=\frac{1}{2}\nabla_{\theta}^{2}\psi$ 可得
\begin{equation*}
\kappa = \frac{1}{2}(\psi_{11}+\psi_{22}).
\tag{8.150}
\end{equation*}
而\textbf{切变}（shear）则使源的形状发生畸变；如果初始为圆形的源被畸变为椭圆度为 $\gamma$、方位角为 $\phi$ 的椭圆，我们定义切变的两个分量为
\begin{align*}
\gamma_1 &= \gamma\cos(2\phi)\\
\gamma_2 &= \gamma\sin(2\phi),
\tag{8.151}
\end{align*}
于是总切变为 $\gamma=\sqrt{\gamma_1^{2}+\gamma_2^{2}}$。用透镜势表出，这些分量为
\begin{align*}
\gamma_1 &= \frac{1}{2}(\psi_{11}-\psi_{22})\\
\gamma_2 &= \psi_{12}=\psi_{21}.
\tag{8.152}
\end{align*}
反解这些关系以求出 $A$ 的分量，得到
\begin{equation*}
A = \begin{pmatrix} 1-\kappa-\gamma_1 & -\gamma_2\\ -\gamma_2 & 1-\kappa+\gamma_1 \end{pmatrix}.
\tag{8.153}
\end{equation*}
于是我们可以用会聚与切变来表示放大率：
\begin{equation*}
\mu = \frac{1}{(1-\kappa)^{2}-\gamma^{2}}.
\tag{8.154}
\end{equation*}

% ---- 书页 355 / p370 ----
透镜效应的这些特征在观测宇宙学中正变得日益重要。显然令人感兴趣的是所谓“强透镜”（strong lensing）的情形：源位于透镜的 Einstein 半径之内，这时可能形成多重像。通过观测同一源的多个像，我们可以推断透镜质量分布的性质（例如用来搜寻暗物质）；我们还可以利用不同路径上的时间延迟来测量 Hubble 常数，并利用透镜事件的统计频率来约束其他宇宙学参数。不过，透镜效应不必很强也能产生重要影响。“弱透镜”（weak lensing）——源与透镜相距超过一个 Einstein 半径——一般导致少量的放大与切变，若不预先知道源的性质便无法探测。然而，切变效应可以用统计方法探测：考察成千上万个星系的形状（假定它们的取向本质上是随机的）。弱透镜切变导致这些形状出现相关联的畸变，从而可以揭示观者与遥远源之间物质分布的大量信息。

\section{我们的宇宙}

在讨论 FRW 宇宙学的行为时，我们屡次提及与我们身处的宇宙相对应的宇宙学参数的实际取值。现在让我们更系统一些：既讨论我们今天看到的宇宙，也讨论向早期时刻的合理外推。我们的讨论必然简略，这既出于篇幅考虑，也因为宇宙学是一个活跃的研究领域；想要了解当前观点的最新描述，请查阅近期的综述文章。

我们对膨胀率的许多直接测定，依靠的是把光度距离公式 (8.123) 应用于某种其内禀光度被假定已知的天体，这类天体称为\textbf{标准烛光}（standard candles）。（偶尔我们也测量其内禀尺度被假定已知的物体的角直径：标准尺［standard rulers］。）例如，Hubble 常数就是用各种标准烛光测量的，不同方法达成的共识收敛于前面提到的值 $H_0=70\pm 10$ km/sec/Mpc。高红移处对线性 Hubble 定律 (8.109) 的偏离可以提供关于密度参数 $\Omega_{i0}$ 的信息，但前提是我们必须拥有非常明亮、且内禀光度已被精确掌握的天体。Ia 型超新星（Type Ia supernovae）恰好提供了这样的天体；它们被认为是白矮星吸积了足够多的质量、以至于超过 Chandrasekhar 极限之后发生的爆发。由于 Chandrasekhar 极限接近普适，相应的爆发在本质上亮度相同（而且一部分内禀变化实际上可以通过追踪亮度随时间的演化而加以扣除）。正是在红移 $z>0.3$ 处对 Ia 型超新星的测量，提供了宇宙学常数非零的第一个直接证据；这些观测意味着 $\Omega_{\Lambda}$ 实际上大于 $\Omega_{\mathrm{M}}$。回想一下：物质是无压强的，$p_{\mathrm{M}}=0$，而真空能则伴随着负压强，$p_{\Lambda}=-\rho_{\Lambda}$。代入第二个弗里德曼方程 (8.68)，我们发现一个同时含有物质与 $\Lambda$ 的宇宙服从

% ---- 书页 356 / p371 ----
\begin{equation*}
\frac{\ddot{a}}{a} = -\frac{4\pi G}{3}(\rho_{\mathrm{M}} - 2\rho_{\Lambda}).
\tag{8.155}
\end{equation*}
于是，如果 $\rho_{\Lambda}$ 相对于 $\rho_{\mathrm{M}}$ 足够大（正如超新星观测所表明的），我们就可以有 $\ddot{a}>0$：一个加速膨胀的宇宙（其含义如 8.4 节所述）。

物质密度本身的测量方法多种多样，通常包括通过寻找成团物质的引力效应来测量密度 $\rho_{\mathrm{M}}$，然后再外推到大尺度。由于 $\rho_{\mathrm{M}}=(3H^{2}/8\pi G)\Omega_{\mathrm{M}}$，用这种方式得到的限值常常以 $\Omega_{\mathrm{M}}h^{2}$ 的形式给出，其中 $h$ 定义于式 (8.70)。如今 $H_0$ 的不确定度看来已足够小，取 $h^{2}\approx 0.5$ 相当稳妥，从此以后我们就这么做。当代的大多数方法都与如下结果一致：
\begin{equation*}
\Omega_{\mathrm{M}0} = 0.3\pm 0.1.
\tag{8.156}
\end{equation*}
在宇宙学常数获得充分的观测证据之前，如此低的物质密度有时曾被当作空间负弯曲（$\kappa<0$）的迹象。

除物质与宇宙学常数之外，宇宙中还有辐射。普通光子是辐射密度中最显然的成分，但任何相对论性粒子都会作出贡献。就光子而言，绝大部分能量密度蕴藏于宇宙微波背景——大爆炸遗留下来的辐射——之中。除光子外，辐射成分唯一显然的候选者是中微子。我们预期遗迹背景中微子的数密度与光子相当；光子的数密度可能还要稍大一些，因为在中微子数目固定下来之后，光子仍能继续产生。然而，如果中微子的质量足够大（大于约 $10^{-4}$ eV），它们今天将已经成为非相对论性的，从而对物质而非辐射作出贡献。目前关于中微子质量的看法提示情况很可能如此，但也并非完全清楚。此外，可以设想在我们已知的粒子之外，还存在迄今尚未探测到的无质量粒子（尽管它们不能太多，否则会抑制大尺度结构的形成）。总而言之，总辐射密度看来很可能与光子密度处于同一数量级；在这种情形下我们将有
\begin{equation*}
\Omega_{\mathrm{R}0} \sim 10^{-4}.
\tag{8.157}
\end{equation*}
如前所述，辐射密度低于物质密度并不令人惊讶，因为随着宇宙的膨胀，前者衰减得更快。辐射密度按 $a^{-4}$ 变化，而物质中的密度按 $a^{-3}$ 变化；因此物质-辐射相等发生的红移为
\begin{equation*}
z_{\mathrm{eq}} \approx \frac{\Omega_{\mathrm{M}0}}{\Omega_{\mathrm{R}0}} \sim 3\times 10^{3}.
\tag{8.158}
\end{equation*}

% ---- 书页 357 / p372 ----
对宇宙学参数的另一个关键约束来自微波背景温度的各向异性。平均温度为 $T_{\mathrm{CMB}}=2.74\,\mathrm{K}$，但 1992 年 COBE 卫星发现了随地点起伏、幅度为 $\Delta T/T\sim 10^{-5}$ 的涨落。这些各向异性有多种来源，其中包括复合时光子移出势阱带来的引力红移/蓝移（Sachs--Wolfe 效应，在大角尺度上占主导）、最后散射面上的内禀温度涨落（在小角尺度上占主导），以及等离子体运动的多普勒效应。描写 CMB 各向异性演化的物理超出了本书的范围。整幅天空的 CMB 温度图显然包含大量信息，但没有任何理论预言任意给定点处的温度应该是多少；现代理论预言的，一般是任意给定角尺度上各向异性大小的期望值。因此，我们把各向异性场按球谐函数展开：
\begin{equation*}
\frac{\Delta T}{T}(\theta,\phi) = \sum_{lm} a_{lm}Y_{lm}(\theta,\phi).
\tag{8.159}
\end{equation*}
$|a_{lm}|^{2}$ 的期望值很可能与 $m$ 无关；否则各向异性的统计特性会随天空中的位置而改变（尽管我们对此应当保持开放的心态）。因此，有待测量的相关参数是
\begin{equation*}
C_l = \langle|a_{lm}|^{2}\rangle.
\tag{8.160}
\end{equation*}
由于对任意固定的 $l$，$m$ 有 $2l+1$ 个可能的取值（从 $-l$ 到 $l$），除了最低的几个 $l$ 之外，对 $a_{lm}$ 的独立测量足以精确确定它们的期望值。在非常小的 $l$ 处这种不可约的不确定度被称为宇宙方差（cosmic variance）。

已有大量实验测量了 $C_l$（所谓的 CMB 功率谱），而改进这些测量在今后若干年内很可能仍是一项重要任务。（除温度各向异性之外，CMB 的偏振也包含大量信息，是另一个需要相当多实验努力的目标。）要把这些观测转化为有用的信息，我们需要一个具体的理论来预言 CMB 功率谱如何作为宇宙学参数的函数而变化。有两种领先的可能性（尽管其中一种比另一种要领先得多）：要么密度扰动在极早的时刻就被印记在所有尺度上——甚至包括物理波长 $\lambda$ 远大于 Hubble 半径 $H^{-1}$ 的模式；要么局部的动力学机制在所有时期充当各向异性的源。后一种可能性已基本上被 CMB 数据排除：如果各向异性是连续产生的，我们预期会看到相当平滑、无特征的 $C_l$ 谱，而观测却显示出显著的结构。因此，更流行的做法是设想某种原初的扰动源，例如暴胀（下一节讨论）。暴胀扰动是绝热的——物质密度中的扰动与辐射密度中的扰动相关联——并且在所有尺度上幅度几乎相等。有了这个输入，我们就可以对 $C_l$ 作出作为所有宇宙学参数函数的明确预言。迄今从实验得到的最重要约束或许是：宇宙在空间上是平直的，或至少非常接近平直；$|\Omega_{\mathrm{c}0}|<0.1$。结合物质密度的测量 $\Omega_{\mathrm{M}}\approx 0.3$，我们得出结论：真空能密度参数应为

% ---- 书页 358 / p373 ----
\begin{equation*}
\Omega_{\Lambda 0} = 0.7\pm 0.1.
\tag{8.161}
\end{equation*}
这与上面描述的 Ia 型超新星结果符合得很好；这里描述的调和图像（concordance picture）正是图 8.4 所指示的那一种。利用 $H_0=70$ km/sec/Mpc 把密度参数换算成物理能量密度，得到
\begin{equation*}
\rho_{\mathrm{vac}} \approx 10^{-8}\ \mathrm{erg/cm^{3}},
\tag{8.162}
\end{equation*}
如 4.5 节讨论真空能时所述。

我们的现今宇宙概貌还有一个值得注意的特征。前面提到，我们宇宙中约 $30\%$ 的能量密度由物质组成。但对宇宙学家而言，“物质”是任何非相对论性粒子的集合；我们通过其引力影响推断出的物质，未必是我们在地球上的经验所熟悉的那种普通物质。\textbf{普通物质}（ordinary matter）指由原子及其成分（质子、中子和电子）构成的一切东西；这包括宇宙中所有的恒星、行星、气体和尘埃，无论能否直接看到。这类物质偶尔也被称为\emph{重子物质}（baryonic matter），其中重子包括质子、中子以及相关的粒子（携带一个被称为重子数的守恒量子数的强相互作用粒子）。当然，电子在概念上是普通物质的重要部分，但按质量计算，与质子和中子相比可以忽略：
\begin{gather*}
m_p = 0.938\ \mathrm{GeV}\\
m_n = 0.940\ \mathrm{GeV}\\
m_e = 0.511\times 10^{-3}\ \mathrm{GeV}.
\tag{8.163}
\end{gather*}
换句话说，普通物质的质量绝大部分来自重子。

事实证明，普通的重子物质远不足以解释观测到的密度 $\Omega_{\mathrm{M}}\approx 0.3$。我们目前对重子密度的最佳估计给出
\begin{equation*}
\Omega_{\mathrm{b}} = 0.04\pm 0.02,
\tag{8.164}
\end{equation*}
这些误差限按大多数标准衡量都是保守的。这一测定来自多种方法：直接计数重子（最不精确的方法）、与 CMB 功率谱（上文讨论过）的一致性，以及与大爆炸核合成（下文讨论）对轻元素丰度的预言相符。因此，物质密度的大部分必定采取\textbf{非重子暗物质}（nonbaryonic dark matter）的形式，我们将把它简称为“暗物质”。（重子也可以是暗的，但越来越普遍的用法是把这一术语保留给非重子成分。）粒子物理标准模型中的每一种已知粒子，作为这种暗物质的候选者都已被排除。幸运的是，标准模型之外还有若干看似合理的候选者，包括中性伴随子（neutralino，超对称性所预言的附加稳定粒子中最轻者，质量 $\geq 100$ GeV）和轴子（axion，因假设的 Peccei--Quinn 对称性自发破缺而产生的轻赝标量粒子，引入该对称性是为了解释强相互作用中 CP 的守恒，质量 $\sim 10^{-4}$ eV）。关于暗物质，我们仅有的几点确切认识之一是：它必须是冷的——不仅今天是非相对论性的，而且在很长一段时间里都必须如此。如果暗物质是热的，它就会自由流出高密度区，抑制星系的形成。关于冷暗物质（cold dark matter，CDM）我们知道的另一件事是，它与普通物质的相互作用必须非常微弱，这样才能解释它为何迄今未被探测到。尽管如此，环境中的暗物质粒子偶尔仍可能与地球实验室里经过仔细屏蔽的探测器发生散射；通过搜寻这种散射的效应来直接探测暗物质，将是今后若干年里另一项重要的实验努力。

% ---- 书页 359 / p374 ----
$\Omega_{\mathrm{M}}=0.3$、$\Omega_{\Lambda}=0.7$ 的这幅图像，似乎能拟合种类多得令人印象深刻的观测数据。这幅图中最令人惊讶的部分是宇宙学常数。在第 4 章中我们提到，对真空能的朴素估计给出的结果比测量值大许多个数量级。实际上有三个相互关联的谜题：为什么宇宙学常数比我们预期的小这么多？构成当前宇宙 $70\%$ 的那份小小的非零能量，其起源是什么？还有，为什么真空能的当前值与物质密度处于同一数量级？最后一个问题尤其严峻，因为真空能与物质密度彼此之间演化得很快：
\begin{equation*}
\frac{\Omega_{\Lambda}}{\Omega_{\mathrm{M}}} \propto a^{3}.
\tag{8.165}
\end{equation*}
如果 $\Omega_{\mathrm{M}}$ 和 $\Omega_{\Lambda}$ 今天彼此相当，那么过去真空能会小得探测不到，而将来物质密度将变得可以忽略。这个“巧合问题”（coincidence problem）迄今仍完全是个谜。有人提出的解法涉及“人择原理”（anthropic principle）：如果宇宙存在许多不同的部分（在空间上，甚至在波函数的不同分支上），其中宇宙学常数取非常不同的数值，那么智慧生命最有可能出现在那些绝对值不太大的地方——一个大的正 $\Lambda$ 会在星系形成之前把粒子撕开，而一个大的负 $\Lambda$ 会使宇宙在生命演化出来之前就再度坍缩。对观测到的真空能的人择解释与数据拟合得很好，尽管为解释这一个量就需要求助如此精巧的方案，在一些人看来有点夸张。

另一种可能与巧合问题有关（也可能无关）的想法是：我们探测到的并非非零的宇宙学常数，而是一个紧密切拟真空能性质的动力学成分。对这种可能性的考虑，促使宇宙学家创造了\textbf{暗能量}（dark energy）这个术语，用来称呼无论实际上是什么的被探测到的东西——不管它是动力学的，还是最终被证明就是宇宙学常数。关于暗能量我们知道的是：它在空间中的分布相对平滑（否则它就会像暗物质那样，通过其局部引力场被探测到），并且随时间演化得很慢（否则它就不会像超新星数据所表明的那样使宇宙加速）。暗能量的一个简单的动力学候选者由缓慢滚动的标量场提供。考虑一个具有通常作用量的场 $\phi$：

% ---- 书页 360 / p375 ----
\begin{equation*}
S = \int \mathrm{d}^{4}x\sqrt{-g}\left[-\frac{1}{2}g^{\mu\nu}\nabla_{\mu}\phi\nabla_{\nu}\phi - V(\phi)\right],
\tag{8.166}
\end{equation*}
其能量-动量张量为
\begin{equation*}
T_{\mu\nu} = \nabla_{\mu}\phi\nabla_{\nu}\phi + \left[\frac{1}{2}g^{\rho\sigma}\nabla_{\rho}\phi\nabla_{\sigma}\phi - V(\phi)\right]g_{\mu\nu}
\tag{8.167}
\end{equation*}
而运动方程为
\begin{equation*}
\Box\phi - \frac{\mathrm{d}V}{\mathrm{d}\phi} = 0.
\tag{8.168}
\end{equation*}
假设场在空间中完全均匀（$\partial_i\phi=0$）。于是利用 Christoffel 符号 (8.44)，我们可以用时间导数和 Hubble 常数表示 d'Alembert 算符，把 (8.168) 写成
\begin{equation*}
\ddot{\phi} + 3H\dot{\phi} + \frac{\mathrm{d}V}{\mathrm{d}\phi} = 0.
\tag{8.169}
\end{equation*}
我们看到，Hubble 参量起到摩擦项的作用；场会趋于沿势滚落，但当 $H$ 太大时运动将被阻尼。因此，具有足够平坦的势的标量场（如图 8.6 所绘）将滚动得非常缓慢，导致动能远小于势能 $V(\phi)$。此时能量-动量张量为
\begin{equation*}
T_{\mu\nu} \approx -V(\phi)g_{\mu\nu},
\tag{8.170}
\end{equation*}
\begin{figure}[H]
\centering
\includegraphics[width=0.72\textwidth]{fig_8.6.pdf}
\caption*{图 8.6\quad 缓慢滚动的标量场的势能。}
\end{figure}
其中 $\phi\approx$ 常数。与式 (4.96) 比较，我们看到标量场势正在摹拟一种真空能。举一个简单的例子，考虑二次势 $V(\phi)=\frac{1}{2}m^{2}\phi^{2}$。这时 (8.169) 描写一个阻尼谐振子，若 $H>m$ 就会出现过阻尼。但在粒子物理单位制中，今天的 Hubble 常数为 $H_0\approx 10^{-33}$ eV，所以这个标量场的质量将不得不比式 (8.163) 中我们熟悉的基本粒子的质量小得难以置信。这看起来像是不自然的微调。尽管如此，动力学暗能量模型正在被积极探索，部分是出于这样的希望：它们能以某种方式通向巧合问题的解决。

% ---- 书页 361 / p376 ----
有了对当代状况的这一认识，我们就可以设想，早期宇宙必须是什么样子，才能产生出我们今天看到的一切。为了物理直观，标记所考虑的时代时，用温度来表示往往比用红移或大爆炸以来的时间更有帮助。今天的温度是
\begin{equation*}
T_0 = 2.74\ \mathrm{K} = 2.4\times 10^{-4}\ \mathrm{eV}.
\tag{8.171}
\end{equation*}
当然，这里说“温度”，我们指的是宇宙微波背景的表观黑体温度；实际上 CMB 自复合以来就不再处于热平衡，所以对这个概念不宜过分当真。在绝热膨胀下，随着每种相对论性粒子红移，温度都在下降，我们有 $T\propto a^{-1}$。但在早期宇宙的特定时刻会发生非绝热相变；在这种情形下温度实际上并不会升高，只是下降得更加缓慢。为了帮助建立温度、密度与尺度因子之间的关系，我们引入\textbf{相对论性自由度的有效数目}（effective number of relativistic degrees of freedom）的两种不同量度：$g_{*}$ 与 $g_{*S}$（其中 $S$ 代表熵）。考虑一组玻色子与费米子物种，每一种都有自己的有效温度 $T_i$ 和自旋态数目 $g_i$。例如，无质量光子有两个自旋态，故 $g_{\gamma}=2$；有质量的 $\frac{1}{2}$ 自旋费米子也有两个自旋态，故 $g_{e^-}=g_{e^+}=2$。相对论性自由度有效数目的这两个不同版本满足
\begin{equation*}
g_{*} = \sum_{\mathrm{bosons}} g_i\left(\frac{T_i}{T}\right)^{4} + \frac{7}{8}\sum_{\mathrm{fermions}} g_i\left(\frac{T_i}{T}\right)^{4}
\tag{8.172}
\end{equation*}
而
\begin{equation*}
g_{*S} = \sum_{\mathrm{bosons}} g_i\left(\frac{T_i}{T}\right)^{3} + \frac{7}{8}\sum_{\mathrm{fermions}} g_i\left(\frac{T_i}{T}\right)^{3}.
\tag{8.173}
\end{equation*}
神秘的 $\frac{7}{8}$ 因子，来源于计算平衡分布函数时玻色统计与费米统计的差别。对任何处于热平衡的物种，其温度 $T_i$ 将等于背景温度 $T$；但也可能有物种在更低的温度下就已退耦，它们对相对论性自由度有效数目的贡献较小。之所以需要定义两种不同的量度，是因为它们扮演的角色不同：第一个通过
\begin{equation*}
\rho_{\mathrm{R}} = \frac{\pi^{2}}{30}g_{*}T^{4},
\tag{8.174}
\end{equation*}
把温度与（相对论性物种的）能量密度联系起来，
而第二个则把温度与尺度因子联系起来：
\begin{equation*}
T \propto g_{*S}^{-1/3}a^{-1}.
\tag{8.175}
\end{equation*}

% ==== tail：8.7 节（Our Universe）的收尾文字，承接书页 361 末的式 (8.175) 之后的段落，
% ==== 至本文件第一个 \section（8.8 暴胀）之前；由装配器移至 ch8_c 之后 ====

% ---- 书页 362 / p377 ----
事实上，只要相对论性自由度还是粒子物理标准模型中的那些，$g_{*}$ 与 $g_{*S}$ 就应当近似相等。一个十分粗略的估算指南是
\begin{equation*}
g_{*} \approx g_{*S} \sim
\left\{
\begin{array}{ll}
100 & T > 300~\mathrm{MeV} \\
10 & 300~\mathrm{MeV} > T > 1~\mathrm{MeV} \\
3 & T < 1~\mathrm{MeV}.
\end{array}
\right.
\tag{8.176}
\end{equation*}
正如我们马上就要讨论的，改变相对论性自由度有效数目的事件，是 $300~\mathrm{MeV}$ 处的 QCD 相变，以及 $1~\mathrm{MeV}$ 处电子/正电子对的湮灭。

有了这些背景，让我们来考虑宇宙从早期到今天的演化。作为开始，我们想象一个 Robertson--Walker 度规，其物质场处于温度 $1~\mathrm{TeV}=1000~\mathrm{GeV}$ 的热平衡中。高温等离子体是基本粒子（夸克、轻子、规范玻色子和 Higgs 玻色子）的复杂混合物。能量密度的主导形式将是相对论性粒子，因此早期宇宙是辐射主导的。它还非常接近平直，因为弗里德曼方程中的曲率项比物质密度和辐射密度演化得更慢。于是弗里德曼方程为
\begin{equation*}
\begin{aligned}
H^{2} &= \frac{8\pi G}{3}\rho_{\mathrm{R}} \\
&\approx 0.1\,g_{*}\,\frac{T^{4}}{\bar{m}_{\mathrm{P}}^{2}},
\end{aligned}
\tag{8.177}
\end{equation*}
其中约化 Planck 标度为 $\bar{m}_{\mathrm{P}}=(8\pi G)^{-1/2}\approx 10^{18}~\mathrm{GeV}$。如果辐射主导的阶段一直延伸到极早期，宇宙的年龄将近似为 $t\sim H^{-1}$，即
\begin{equation*}
t \sim \frac{\bar{m}_{\mathrm{P}}}{T^{2}}.
\tag{8.178}
\end{equation*}
用常规单位表示，这变为
\begin{equation*}
t \sim 10^{-6}\left(\frac{\mathrm{GeV}}{T}\right)^{2}~\mathrm{sec}.
\tag{8.179}
\end{equation*}

目前粒子加速器上的实验已经为大约 $100~\mathrm{GeV}$ 以下的物理图景提供了准确的描述，因此再外推一个数量级仍属于合理外推的范围。在更高的温度下我们就不太确定会发生什么了；在 $1~\mathrm{TeV}$ 和 Planck 标度之间也许没有任何非常有趣的事情，也可能这个区域充满了各式各样的意外。当然，也可以设想宇宙学在更低的温度下就会给出意外，尽管标准模型物理已被理解得很透彻；本节我们描述的是一个保守的方案，但一如既往，保持开放的头脑总是值得的。

% ---- 书页 363 / p378 ----
标准模型的一个关键特征是电弱部分的自发破缺对称性。在宇宙学中，这种对称性破缺发生在电弱相变处，温度 $T\sim 200~\mathrm{GeV}$。在此温度之上对称性尚未破缺，因此基本费米子（夸克和轻子）与弱相互作用规范玻色子都无质量；在此温度之下，我们便得到低能实验所熟知的那些质量模式。电弱相变预计不会给晚期宇宙留下任何可察觉的影响；一个可能的例外是重子生成（baryogenesis），下面讨论。

在这些温度下，量子色动力学（QCD）所描述的强相互作用并不那么强。在低能量/温度下，QCD 表现出“禁闭”（confinement）——夸克和胶子被束缚成重子和介子这样的复合粒子。但在 QCD 标度 $\Lambda_{\mathrm{QCD}}\sim 300~\mathrm{MeV}$ 之上，夸克和胶子是自由粒子。随着宇宙膨胀和冷却，强相互作用粒子被禁闭成束缚态，这正是 (8.176) 中所指出的相对论性自由度有效数目第一次下降的原因。QCD 相变预计不会在可观测宇宙上留下显著的印记。

正如强相互作用在高温下并不很强，弱相互作用也不像你可能以为的那么弱；它们仍然是弱的——在微扰论能够准确描述的意义上——但其发生得足够快，足以使中微子这类弱相互作用粒子保持在热平衡中。当 $T\sim 1~\mathrm{MeV}$ 时，情况不再如此。这也大致是电子和正电子变为非相对论性并发生湮灭的温度，相对论性自由度的有效数目因此减少，但这两件事互不相干。对 $1~\mathrm{MeV}$ 以下的温度，我们说弱相互作用已经“冻结”（frozen out）——相互作用速率降到宇宙的膨胀速率之下，相互作用发生得过于稀少，无法再使粒子保持平衡。冷暗物质粒子有可能就是在这个温度从等离子体中退耦的。更有把握的是，我们可以断定中子和质子不再相互转换。在这个温度下中子的平衡丰度约为质子丰度的 $1/6$（因为中子的质量稍大）。中子的寿命有限（$\tau_{n}=890~\mathrm{sec}$），只比这个时期宇宙的年龄 $t(1~\mathrm{MeV})\approx 1~\mathrm{sec}$ 略长，但它们已开始逐渐衰变为质子和轻子。然而不久之后，我们就到达略低于 $100~\mathrm{keV}$ 的温度，\textbf{大爆炸核合成}（Big-Bang Nucleosynthesis，BBN）开始了。

每个核子的核结合能典型地为 $1~\mathrm{MeV}$ 的量级，所以你也许会以为核合成应该更早发生；然而，每个核子对应的光子数太多，阻碍了核合成的发生，直到温度降到 $100~\mathrm{keV}$ 以下。在那一点上，中子/质子比约为 $1/7$。在所有轻核之中，核子驻留在 $^{4}\mathrm{He}$ 里在能量上是有利的，事实上大多数自由中子也确实被转化成了 $^{4}\mathrm{He}$；每两个中子和十四个质子，最终得到一个氦核和十二个质子。于是，按质量计约有 $25\%$ 的重子被转化为氦。此外，还有痕量的氘（每质子约 $10^{-5}$ 个氘核）、$^{3}\mathrm{He}$（同样约 $10^{-5}$）以及 $^{7}\mathrm{Li}$（约 $10^{-10}$）。

% ---- 书页 364 / p379 ----
当然，这些数字是预言，而它们确实被对轻元素原初丰度的观测所证实。（更重的元素并非在大爆炸中合成，而是需要晚期宇宙中的恒星过程。）我们对许许多多关键细节作了粗略带过，尤其是那些解释各种丰度如何依赖宇宙学参数的细节。例如，设想我们偏离标准模型，引入三种以上的轻中微子物种。这会通过式 (8.174) 在固定温度下增大辐射能量密度，转而又缩短了与给定温度相联系的时标（因为 $t\sim H^{-1}\propto\rho_{\mathrm{R}}^{-1/2}$）。核合成因此会稍早发生，导致中子丰度更高，从而 $^{4}\mathrm{He}$ 丰度也更大。对原初氦丰度的观测与标准模型的预言相符，给出了轻中微子数目接近三的第一个证据。类似地，与核合成相联系的所有温度和时标都依赖于重子/光子比；要与观测到的丰度相符，就要求每光子大约有 $5\times 10^{-10}$ 个重子，这正是重子密度参数之估计值 (8.164) 的由来，以及随之而来的对非重子暗物质的需要。

就我们当下的目的而言，原初核合成最深刻的特征，也许是它对温度与膨胀率之间的弗里德曼关系——从而对 Einstein 方程——的敏感依赖。BBN 的成功为 GR 提供了一个在远离我们日常经验的领域中的严格检验。Einstein 的理论当初被提出，主要是出于把引力与电磁学的洛伦兹对称性不变性相协调的需要；而这样一个理论竟能成功描述宇宙仅一秒钟大小时的增长，这实在是令人赞叹的成就。时至今日，BBN 仍是对各种替代引力理论最强有力的约束之一；特别地，它是我们拥有任何直接观测特征的最早时期。

核合成之后，我们得到一个由质子、电子和光子主导的等离子体，其中还有一些氦和其他核。也有暗物质存在，但假定到这个时期它已不与普通物质相互作用。下一个重要事件直到\textbf{复合}（recombination）才发生：电子与质子结合（它们与氦的结合稍早一些）。复合发生在温度 $T\approx 0.3~\mathrm{eV}$；此时宇宙是物质主导的。同样地，由于氢的结合能是 $13.6~\mathrm{eV}$，你也许会以为复合应该更早发生，但巨大的光子/重子比把它推迟了。复合至关重要的意义在于，它标志着宇宙变得透明的时刻。环境光子与自由电子强烈地相互作用，因此在复合之前光子的平均自由程非常短；而一旦电子和质子结合成中性氢，平均自由程实际上就变成了无穷大。这些环境光子今天作为宇宙微波背景被我们看到，它提供了宇宙在 $T\approx 0.3~\mathrm{eV}$ 时、即红移 $z\approx 1200$ 时的一幅快照。复合是一个颇为渐进的过程，因此对它发生于何时所作的任何指定都必然是近似的。

% ---- 书页 365 / p380 ----
复合之后，宇宙要走过一段被称为“黑暗时代”（dark ages）的漫长时期，其间星系通过引力不稳定性逐渐集结起来，但还没有可见的恒星来照亮宇宙。黑暗时代是一个神秘的时期；恒星和星系形成的过程高度复杂且非线性，无疑需要新类型的观测，我们才能对这一时代有充分的理解。

我们的故事现在已经把我们带到了今天，但还有几个缺失的环节需要回过头来补上。其一是宇宙中物质与反物质之间的不对称。宇宙中基本上全部可见物质似乎都由质子、中子和电子构成，而非它们的反粒子；如果遥远的星系主要由反物质构成，我们就应当能观测到物质/反物质区域边界处质子与反质子偶尔湮灭所产生的高能光子。虽然可以把这种不对称作为初始条件硬性加进去，但这总显得有些不能令人满意，大多数物理学家更愿意找到某种重子生成的动力学机制，让一个初始时物质/反物质对称的态演化成我们今天的宇宙。这样的破缺对称在粒子物理中很常见，事实上人们已经提出了许多重子生成机制（通常在电弱标度或更高的温度下）。然而，这些具体方案中没有一个已被证明足够令人信服而被采纳为标准方案。要理解重子不对称的起源，我们很可能需要对超越标准模型的物理有更好的理解。

另一个需要提及的缺失特征是，宇宙当然并非完全均匀和各向同性；宇宙中目前的大尺度结构似乎演化自极早期就存在的、幅度为 $\delta\rho/\rho\sim 10^{-5}$ 的绝热且近无标度微扰。这些微扰的绝热性和近无标度性的证据，来自对 CMB 和大尺度结构的联合观测。高度的各向同性和均匀性，以及对其的小偏离，在传统宇宙学中都只是作为神秘的初始条件被强加上去的。暴胀宇宙方案（inflationary-universe scenario）为两者提供了一个可能的动力学起源，我们现在就转向它。

\section{暴胀}

在对大爆炸模型的传统理解中，宇宙被认为在早期是辐射主导的，在晚期是物质主导的，而且——正如我们现在怀疑的——在非常晚近的时候才转变为真空主导。这幅图像在描述形形色色的观测数据方面取得了巨大成功；尽管如此，我们仍然可以问，产生这样一个宇宙的初始条件是否显得自然。这是那种人们在宇宙学中会问、在其他科学中却不会问的问题。通常，作为物理学家，我们寻找自然定律，并想象我们可以自由地指定初始条件，然后问它们在这些定律下如何演化。但宇宙似乎只有一组初始条件，因此，想知道它们究竟是相对平常的还是经过精细调节的，似乎是合情合理的。

% ---- 书页 366 / p381 ----
在传统图像中，早期宇宙确实被精细调节到了难以置信的精度。特别是，我们宇宙的两个特征看上去非常不寻常：它的空间平直性，以及它高度的各向同性和均匀性。有可能这就是我们恰好被给予的宇宙，追问不同初始条件的可能性有多大毫无意义；另一种可能是，如果存在某种动力学机制，能把宽广谱系的初始条件朝平直性和均匀性/各向同性的方向演化，那么这些条件就比乍看之下更有可能出现。暴胀宇宙方案提供了这样一种机制（而且还不止如此），它已成为现代宇宙学的核心组织原理，即使我们还远不能证明它的正确性。

在描述暴胀之前，让我们先描述它声称解决的两个不自然性问题：与均匀性/各向同性相联系的\textbf{平坦性问题}（flatness problem）和\textbf{视界问题}（horizon problem）。平坦性问题来自考虑一个有物质和辐射但没有真空能的宇宙中的弗里德曼方程；为方便后文，我们用约化 Planck 质量 $\bar{m}_{\mathrm{P}}=(8\pi G)^{-1/2}$ 把它写成
\begin{equation*}
H^{2} = \frac{1}{3\bar{m}_{\mathrm{P}}^{2}}\left(\rho_{\mathrm{M}}+\rho_{\mathrm{R}}\right) - \frac{\kappa}{a^{2}}.
\tag{8.180}
\end{equation*}
曲率项 $-\kappa/a^{2}$ 与 $a^{-2}$ 成正比（这显而易见），而能量密度项随尺度因子的增大衰减得更快，$\rho_{\mathrm{M}}\propto a^{-3}$、$\rho_{\mathrm{R}}\propto a^{-4}$。这就引出了一个问题：考虑到自 Planck 时期以来 $a$ 已经增大了大约 $10^{30}$ 倍，为什么比值 $(\kappa a^{-2})/(\rho/3\bar{m}_{\mathrm{P}}^{2})$ 没有远远大于 1？换句话说，在物质/辐射主导的宇宙中，$\Omega=1$ 这个点是一个排斥性的不动点——对这个值的任何偏离都会随时间增长，那么为什么我们今天观测到的是 $\Omega\approx 1$？

视界问题源于 FRW 宇宙学中粒子视界的存在，如图 8.7 所示。视界之所以存在，是因为自大爆炸奇点以来只有有限的时间，因而在宇宙的年龄之内光子只能走过有限的距离，正如我们在第 2 章中简要讨论过的。考虑在平直宇宙中沿径向轨迹运动的一个光子（推广到非平直宇宙是直截了当的）。径向类光路径满足
\begin{equation*}
0 = \mathrm{d}s^{2} = -\mathrm{d}t^{2} + a^{2}\mathrm{d}r^{2},
\tag{8.181}
\end{equation*}
因此这样的一个光子在 $t_{1}$ 与 $t_{2}$ 两时刻之间走过的共动（坐标）距离为
\begin{equation*}
\Delta r = \int_{t_{1}}^{t_{2}}\frac{\mathrm{d}t}{a(t)}.
\tag{8.182}
\end{equation*}
要得到任一时刻 $t$ 处观者所测得的物理距离，只需乘上 $a(t)$。为简单起见，让我们想象自己处在一个物质主导的
% ---- 书页 367 / p382 ----
宇宙，对它而言
\begin{equation*}
a = \left(\frac{t}{t_{0}}\right)^{2/3}.
\tag{8.183}
\end{equation*}
\begin{figure}[H]
\centering
\includegraphics[width=0.72\textwidth]{fig_8.7.pdf}
\caption*{图 8.7\quad 从一个由大爆炸奇点膨胀而来的宇宙中的过去光锥，展示了宇宙学中的粒子视界。处于复合时刻的一些点——今天作为天空中相反两侧的宇宙微波背景被我们观测到——其过去光锥互不重叠（在传统宇宙学中）；没有任何因果信号能够影响它们，使它们具有相同的温度。}
\end{figure}
记住 $a_{0}=1$。因此 Hubble 参量为
\begin{equation*}
\begin{aligned}
H &= \frac{2}{3}t^{-1} \\
&= a^{-3/2}H_{0}.
\end{aligned}
\tag{8.184}
\end{equation*}
于是该光子走过的共动距离为
\begin{equation*}
\Delta r = 2H_{0}^{-1}\left(\sqrt{a_{2}}-\sqrt{a_{1}}\right).
\tag{8.185}
\end{equation*}
在尺度因子取任一固定值 $a=a_{*}$ 时的共动视界尺度，是光子自大爆炸以来走过的距离，
\begin{equation*}
r_{\mathrm{hor}}(a_{*}) = 2H_{0}^{-1}\sqrt{a_{*}}.
\tag{8.186}
\end{equation*}
而在 $a_{*}$ 处的空间超曲面上测得的物理视界尺寸，因此就简单地是
\begin{equation*}
d_{\mathrm{hor}}(a_{*}) = a_{*}r_{\mathrm{hor}}(a_{*}) = 2H_{*}^{-1}.
\tag{8.187}
\end{equation*}
事实上，对任何包含物质和辐射混合物的近平直宇宙，在任一给定时期我们都有
\begin{equation*}
d_{\mathrm{hor}}(a_{*}) \sim H_{*}^{-1} = d_{H}(a_{*}),
\tag{8.188}
\end{equation*}

% ---- 书页 368 / p383 ----
其中 Hubble 距离 $d_{H}$ 已在式 (8.71) 中引入。这种近似相等带来一种强烈的诱惑，即把“视界距离”和“Hubble 距离”当作同义词使用；这种诱惑应当抵制，因为暴胀可以使前者远大于后者，我们很快就会证明这一点。

视界问题简单说来就是：CMB 在很高的精度上是各向同性的，尽管最后散射面上相距遥远的各点彼此完全处于对方的视界之外。当我们观看 CMB 时，我们观测的是尺度因子 $a_{\mathrm{CMB}}\approx 1/1200$ 处的宇宙；由式 (8.185)，CMB 上一点与地球上的观者之间的共动距离为
\begin{equation*}
\begin{aligned}
\Delta r &= 2H_{0}^{-1}\left(1-\sqrt{a_{\mathrm{CMB}}}\right) \\
&\approx 2H_{0}^{-1}.
\end{aligned}
\tag{8.189}
\end{equation*}
然而，这样一个点所对应的共动视界距离为
\begin{equation*}
\begin{aligned}
r_{\mathrm{hor}}(a_{\mathrm{CMB}}) &= 2H_{0}^{-1}\sqrt{a_{\mathrm{CMB}}} \\
&\approx 6\times 10^{-2}H_{0}^{-1}.
\end{aligned}
\tag{8.190}
\end{equation*}
因此，如果我们观测 CMB 中相距遥远的两个部分，它们的视界将互不重叠；CMB 天空中不同的斑块在复合时刻是因果不连通的。然而，观测却表明它们具有相同的温度，且精度很高。于是问题就来了：既然它们从未处于因果接触之中，它们是如何预先知道要以正确的方式协调各自的演化的？我们必须设法修正传统 FRW 宇宙学的因果结构。

让我们考虑这样修正传统图像：设定一段\textbf{暴胀}（inflation）时期，即极早期宇宙中一段加速膨胀（$\ddot{a}>0$）的时期，由物质或辐射以外的某种成分驱动，这种成分随宇宙膨胀而红移得很慢。这样一来，平坦性问题和视界问题便可以同时解决。为简单起见，考虑暴胀由恒定的真空能驱动、从而导致指数膨胀的情形。于是，在真空主导的时期，$\rho/3\bar{m}_{\mathrm{P}}^{2}\propto a^{0}$ 相对于 $-\kappa/a^{2}$ 迅速增大，宇宙随时间变得越来越平（$\Omega$ 被驱动到 1）。如果这个过程持续了足够长的时间，然后真空能才转化为物质和辐射，那么密度参数将足够接近 1，以至于进入当前时代之前它还来不及发生可察觉的改变。与此同时，视界问题可以追溯到这样一个事实：任意两个共动物体之间的物理距离随尺度因子增长，而物质主导或辐射主导宇宙中的物理视界尺寸增长得更快，$d_{\mathrm{hor}}\sim a^{n/2}H_{0}^{-1}$。这同样可以由早期一段指数膨胀时期来解决：真实的视界尺寸会增长到一个惊人的数量，以至于我们今天的视界实际上远大于“视界等于 Hubble 半径 $H_{0}^{-1}$”这个天真的估计。

事实上，真正的指数膨胀并非必需；对任何加速膨胀，空间曲率相对于能量密度都会减小，
% ---- 书页 369 / p384 ----
而视界距离将迅速增长。通常我们要求这个加速时期持续 60 个或更多的 $e$ 指数倍（e-fold，$e$ 指数倍数为 $N=\Delta\ln a$），这正是解决视界问题所需要的。很容易做过头，暴胀一般会使今天的宇宙在空间上平直得难以置信。

现在让我们考虑如何在早期宇宙中得到一个暴胀阶段。最直接的办法是利用一个标量场的势所提供的真空能，这个场就是\textbf{暴胀子}（inflaton）。想象一个由空间均匀的标量（场）的能量主导的宇宙。相关的运动方程正是我们在 8.7 节中讨论动力学暗能量时的那些方程；唯一的区别是暴胀的能量标度要高得多。我们有 RW 度规中标量场的运动方程，
\begin{equation*}
\ddot{\phi} + 3H\dot{\phi} + V'(\phi) = 0,
\tag{8.191}
\end{equation*}
以及弗里德曼方程
\begin{equation*}
H^{2} = \frac{1}{3\bar{m}_{\mathrm{P}}^{2}}\left(\frac{1}{2}\dot{\phi}^{2} + V(\phi)\right).
\tag{8.192}
\end{equation*}
我们忽略了曲率项，因为暴胀反正会把宇宙拉平。如果场的演化足够平缓，使势能得以主导动能，并且 $\phi$ 的二阶导数足够小，足以让这种状态维持一段充分长的时间，暴胀就能发生。因此，我们希望
\begin{equation*}
\begin{aligned}
\dot{\phi}^{2} &\ll V(\phi), \\
|\ddot{\phi}| &\ll |3H\dot{\phi}|,\ |V'|.
\end{aligned}
\tag{8.193}
\end{equation*}
满足这些条件要求两个无量纲量足够小，它们就是所谓的\textbf{缓慢滚动参数}（slow-roll parameters）：
\begin{equation*}
\begin{aligned}
\epsilon &= \frac{1}{2}\bar{m}_{\mathrm{P}}^{2}\left(\frac{V'}{V}\right)^{2}, \\
\eta &= \bar{m}_{\mathrm{P}}^{2}\left(\frac{V''}{V}\right).
\end{aligned}
\tag{8.194}
\end{equation*}
注意 $\epsilon\ge 0$，而 $\eta$ 可正可负。还要注意，这些定义并非人人通用；有些人喜欢用 Hubble 参量而不是势来定义它们。我们的选择描述的是场是否有机会缓慢滚动一阵子；而用 Hubble 参量的描述给出的则是场实际上是否正在缓慢滚动。当这两个量都很小时，我们就可以有一个长久的暴胀阶段。然而它们并不充分；无论势是什么样子，我们总可以选取初始条件使 $|\dot{\phi}|$ 大到缓慢滚动永不适用。不过，只要缓慢滚动参数很小，大多数初始条件都会被吸引到暴胀阶段去。

% ---- 书页 370 / p385 ----
要发明满足缓慢滚动条件的势并不难。考虑也许是最简单的例子\footnote{我们沿用 A.R.~Liddle, \emph{An Introduction to Cosmological Inflation}（\texttt{http://arxiv.org/astro-ph/9901124}）中的讲述。}，$V(\phi)=\frac{1}{2}m^{2}\phi^{2}$。在这种情形下
\begin{equation*}
\epsilon = \eta = \frac{2\bar{m}_{\mathrm{P}}^{2}}{\phi^{2}}.
\tag{8.195}
\end{equation*}
显然，只要 $\phi$ 足够大，我们就能让缓慢滚动参数要多小有多小。然而，我们有一个约束：能量密度不应高达 Planck 标度，这样我们的经典分析才有意义；这意味着 $\phi\ll\bar{m}_{\mathrm{P}}^{2}/m$。如果我们把场从值 $\phi_{i}$ 开始，那么在暴胀结束之前（也就是缓慢滚动参数变成 1 的量级之前）的 $e$ 指数倍数将是
\begin{equation*}
\begin{aligned}
N &= \int_{t_{i}}^{t_{e}} H\,\mathrm{d}t \\
&\approx -\bar{m}_{\mathrm{P}}^{-2}\int_{\phi_{i}}^{\phi_{e}}\frac{V}{V'}\,\mathrm{d}\phi \\
&\approx \frac{\phi_{i}^{2}}{4\bar{m}_{\mathrm{P}}^{2}} - \frac{1}{2}.
\end{aligned}
\tag{8.196}
\end{equation*}
第一个等式总是成立，第二个等式用了缓慢滚动近似，第三个等式是这个具体模型的结果。要得到 60 个 $e$ 指数倍，我们需要 $\phi_{i}>16\bar{m}_{\mathrm{P}}$。再加上能量密度的上限，我们发现质量参数存在一个上限：$m\ll\bar{m}_{\mathrm{P}}/16$。事实上，观测到的密度涨落的大小对 $m$ 给出了更严格的上限，我们将在下面讨论。但 $m$ 没有下限，因此，只有当我们愿意设定 $m\ll\bar{m}_{\mathrm{P}}$ 这样的巨大层级差，或者等价地一个很小的无量纲数 $m/\bar{m}_{\mathrm{P}}$ 时，才容易得到合适的暴胀势。用 $\lambda\phi^{4}$ 势做同样的练习也会得出类似的结论：$\lambda$ 必须相当小；我们常说暴胀子必须是弱耦合的。当然，从某种意义上说这是在作弊，因为当场值 $\phi>\bar{m}_{\mathrm{P}}$ 时，我们应当预期有效势中会出现形如 $\bar{m}_{\mathrm{P}}^{4-n}\phi^{n}$（$n>4$）的附加项并变得重要。所以在一个\emph{现实}的模型中，要得到合适的势可能相当困难。

暴胀终会在某个时刻结束，暴胀子势中的能量将转化为物质和辐射的热化气体，这个过程被称为“再热”（reheating）。恰当地理解再热过程至关重要，因为它控制着我们宇宙中种种我们想要或不想要的遗留物（relic）的产生。例如，暴胀的一个重要好处是，它能够把早期宇宙中可能产生、但今天并未观测到的各种遗留物“暴胀掉”（inflate away）。一个经典的例子出现在粒子物理大统一理论的语境中：大统一理论一般预言存在超重的磁单极子，其丰度比观测容许的还要多出许多个数量级。

% ---- 书页 371 / p386 ----
历史上，磁单极子问题是 Guth 发明暴胀的首要动机；平坦性问题和视界问题的解决当时被视为额外的红利。暴胀可以把磁单极子的丰度稀释到合适的程度，但如果宇宙再热到超过大统一相变的温度，它们又会重新产生；幸运的是，对大多数模型而言这并不是一个严格的约束。类似的考虑也适用于其他不想要的遗留物；在超对称模型中，引力微子（gravitino，引力子的超对称伙伴）的丰度引起一个尤其令人头疼的问题。与此同时，又必须再热到足够高的温度，以便某种重子生成方案得以进行。对于在某个粒子物理模型内实现暴胀的任何具体方案，关键的检查是：不想要的遗留物被清除，而想要的遗留物（例如重子）得以保留。

暴胀方案的一个关键要素是密度微扰的产生，它们可能是我们今天观测到的 CMB 温度各向异性和星系大尺度结构的起源。暴胀产生密度微扰背后的想法相当直截了当。暴胀会把任何环境粒子密度迅速衰减到零，只留下真空。但是加速宇宙中的真空态具有非零的温度，即 Gibbons--Hawking 温度，它类似于黑洞的 Hawking 温度。我们无法在本书中深入探讨这个主题；这里只概述其基本结果。

对于一个由势能 $V$ 主导的宇宙，Gibbons--Hawking 温度为
\begin{equation*}
T_{\mathrm{GH}} = \frac{H}{2\pi} \sim \frac{V^{1/2}}{\bar{m}_{\mathrm{P}}}.
\tag{8.197}
\end{equation*}
与这个温度相对应的，是暴胀子场 $\phi$ 在每个波数 $k$ 处的涨落，其大小为
\begin{equation*}
|\Delta\phi|_{k} = T_{\mathrm{GH}}.
\tag{8.198}
\end{equation*}
由于势按假设近乎平坦，$\phi$ 的涨落会导致能量密度的小涨落，
\begin{equation*}
\delta\rho = V'(\phi)\delta\phi.
\tag{8.199}
\end{equation*}
因此，暴胀在所有尺度上都产生密度微扰。微扰的幅度在每个波数处都几乎相等，但由于暴胀子滚动时 $V$ 的逐渐变化，会有微小的偏离。描述这些微扰是一个繁乱的课题，涉及数不清的各种记号。一个合理的出发点是方均根（root-mean-square，RMS）密度涨落，
\begin{equation*}
\left.\frac{\delta\rho}{\rho}\right|_{\mathrm{rms}} = \sqrt{\left\langle\left(\frac{\delta\rho}{\rho}\right)^{2}\right\rangle},
\tag{8.200}
\end{equation*}

% ---- 书页 372 / p387 ----
其中尖括号表示对空间位置的平均。对于统计各向同性的微扰（期望幅度与方向无关），稍稍做一点傅里叶分析便能让我们写出
\begin{equation*}
\left(\left.\frac{\delta\rho}{\rho}\right|_{\mathrm{rms}}\right)^{2} = \int \Delta^{2}(k)\,\mathrm{d}(\ln k),
\tag{8.201}
\end{equation*}
其中我们引入了无量纲功率谱，
\begin{equation*}
\Delta^{2}(k) \equiv \frac{k^{3}|\delta_{k}|^{2}}{2\pi^{2}},
\tag{8.202}
\end{equation*}
而 $\delta_{k}$ 是分数密度微扰（fractional density perturbation）的傅里叶变换的期望值，
\begin{equation*}
\delta_{\mathbf{k}} = \frac{1}{(2\pi)^{3/2}}\int e^{-i\mathbf{k}\cdot\mathbf{x}}\frac{\delta\rho}{\rho}\,\mathrm{d}^{3}x,
\tag{8.203}
\end{equation*}
这里我们已假设它是各向同性的。无量纲功率谱是时间的函数，因为每个模式的幅度都在演化；表达任何具体模型的预言时，最常用的做法是取模式的物理波长 $\lambda=a/k$ 等于 Hubble 半径 $H^{-1}$ 的那一时刻的微扰幅度，
\begin{equation*}
A_{S}^{2}(k) \equiv \left.\Delta^{2}(k)\right|_{k=aH}.
\tag{8.204}
\end{equation*}
这样，$A_{S}(k)$ 量度的是不同模式在不同时刻的幅度。对于由缓慢滚动标量场驱动的暴胀，$A_{S}(k)$ 与势通过下式相联系：
\begin{equation*}
A_{S}^{2}(k) \sim \left.\frac{V^{3}}{\bar{m}_{\mathrm{P}}^{6}(V')^{2}}\right|_{k=aH} \sim \left.\frac{V}{\bar{m}_{\mathrm{P}}^{4}\epsilon}\right|_{k=aH}.
\tag{8.205}
\end{equation*}
我们有意略去了无量纲的数值因子——它们在不同文献之间差别很大——以便突出对势的依赖。

谱带着下标“S”，是因为它描述的是度规中的标量涨落。这些涨落与能量-动量分布绑定在一起，暴胀产生的密度涨落是绝热的——所有物种的密度涨落都是相互关联的。涨落还是高斯的，其含义是：描述不同尺度上涨落的那些傅里叶模式的相位互不相关。暴胀微扰的这些方面——近无标度的绝热密度涨落谱、高斯分布——都与当前对 CMB 和大尺度结构的观测一致，而今后几年预定要收集的新数据应会大大提高这些检验的精度。

% ---- 书页 373 / p388 ----
在暴胀期间被激发的不仅是近无质量的暴胀子，还有任何其他近无质量的粒子。另一个重要的例子是引力子，它对应于度规中的张量微扰（引力场的传播激发）。张量涨落具有谱
\begin{equation*}
A_{T}^{2}(k) \sim \left.\frac{V}{\bar{m}_{\mathrm{P}}^{4}}\right|_{k=aH}.
\tag{8.206}
\end{equation*}
重要的是，张量幅度只依赖于势本身，而不依赖于它的导数；因此，对张量微扰的观测将直接给出关于暴胀能量标度的信息。

为了理解观测，用可观测的量来参数化微扰谱是有用的。于是我们写出
\begin{equation*}
A_{S}^{2}(k) \propto k^{n_{S}-1}
\tag{8.207}
\end{equation*}
以及
\begin{equation*}
A_{T}^{2}(k) \propto k^{n_{T}},
\tag{8.208}
\end{equation*}
其中 $n_{S}$ 和 $n_{T}$ 是谱指数（spectral index）。它们与势的缓慢滚动参数的关系为
\begin{equation*}
n_{S} = 1 - 6\epsilon + 2\eta
\tag{8.209}
\end{equation*}
以及
\begin{equation*}
n_{T} = -2\epsilon.
\tag{8.210}
\end{equation*}
在我们所考虑的这类模型（由单个缓慢滚动的标量场驱动）中，存在一个把标量模式和张量模式的幅度与谱指数联系起来的自洽关系。它可以用一种与约定无关的方式表达为可观测量之间的关系，即不同微扰所导致的温度涨落 $\Delta T$ 之间的关系：
\begin{equation*}
\frac{(\Delta T/T)_{T}^{2}}{(\Delta T/T)_{S}^{2}} = -7n_{T}.
\tag{8.211}
\end{equation*}

张量微扰的存在是暴胀的一个关键预言，原则上可以通过对 CMB 偏振的观测加以检验。偏振也可以由通常的密度涨落诱导产生，其机制是非均匀等离子体中 Thompson 散射截面的各向异性。幸运的是，我们可以设想把天空中的偏振矢量场分解为无旋部分（$E$ 模式）和有旋部分（$B$ 模式）；标量微扰导致 $E$ 模式偏振，而张量微扰导致 $B$ 模式（除去复合之后的宇宙中某些不可避免的演化加工）。CMB 偏振已经被探测到了；未来的挑战将在于把标量和张量的贡献分离出来，以检验简单暴胀模型的预言 (8.211)。

% ---- 书页 374 / p389 ----
当然，这不仅要求探测到张量诱导的偏振，还要求以一定的精度测量它的谱指数。

我们对微扰幅度已有的知识，已经给了我们关于暴胀能量标度的重要信息。张量微扰只依赖 $V$ 本身，不依赖它的导数；如果 COBE 观测到的 CMB 各向异性源自张量涨落（有可能，尽管不太可能），我们可以立刻导出 $V_{\mathrm{inflation}}\sim(10^{16}~\mathrm{GeV})^{4}$。这里，受到约束的 $V$ 值是产生观测涨落的那一时刻的值，即暴胀结束前 60 个 $e$ 指数倍时的值。这极其让人联想到大统一标度，十分令人鼓舞。即使在更可能的情形下——CMB 中观测到的微扰本质上是标量的——我们仍然可以写出
\begin{equation*}
V_{\mathrm{inflation}}^{1/4} \sim \epsilon^{1/4}10^{16}~\mathrm{GeV},
\tag{8.212}
\end{equation*}
其中 $\epsilon$ 是式 (8.194) 定义的缓慢滚动参数。虽然我们预期 $\epsilon$ 很小，但指数中的 $1/4$ 意味着对 $\epsilon$ 的依赖相当弱；除非这个参数小得出奇，否则很有可能 $V_{\mathrm{inflation}}^{1/4}\sim 10^{15}$--$10^{16}~\mathrm{GeV}$。我们能对如此惊人的能量标度获得这样的信息，这一事实本身就足以令人惊叹。

\section{习题}

\textbf{1.}\quad 考虑一个 $(N+n+1)$ 维时空，坐标为 $\{t,\,x^{I},\,y^{i}\}$，其中 $I$ 从 1 取到 $N$，$i$ 从 1 取到 $n$。设度规为
\begin{equation*}
\mathrm{d}s^{2} = -\mathrm{d}t^{2} + a^{2}(t)\delta_{IJ}\,\mathrm{d}x^{I}\mathrm{d}x^{J} + b^{2}(t)\gamma_{ij}(y)\,\mathrm{d}y^{i}\mathrm{d}y^{j},
\tag{8.213}
\end{equation*}
其中 $\delta_{IJ}$ 是通常的 Kronecker delta，而 $\gamma_{ij}(y)$ 是 $n$ 维最大对称空间流形上的度规。设想我们把度规 $\gamma$ 归一化，使得曲率参数
\begin{equation*}
k = \frac{R(\gamma)}{n(n-1)}
\tag{8.214}
\end{equation*}
等于 $+1$、$0$ 或 $-1$，其中 $R(\gamma)$ 是度规 $\gamma_{ij}$ 所对应的 Ricci 标量。

\textbf{(a)}\quad 计算这个度规的 Ricci 张量。

\textbf{(b)}\quad 用能量密度 $\rho$ 以及 $x^{I}$ 和 $y^{i}$ 方向上的压强 $p^{(N)}$、$p^{(n)}$ 定义一个能量-动量张量：
\begin{gather*}
T_{00} = \rho \tag{8.215}\\
T_{IJ} = a^{2}p^{(N)}\delta_{IJ} \tag{8.216}\\
T_{ij} = b^{2}p^{(n)}\gamma_{ij}. \tag{8.217}
\end{gather*}
把这个度规和 $T_{\mu\nu}$ 代入 Einstein 方程，推导 $a$ 和 $b$ 所满足的类弗里德曼方程（共三个独立方程）。

% ---- 书页 375 / p390 ----
\textbf{(c)}\quad 推导静态解（即 $\dot{a}=\dot{b}=\ddot{a}=\ddot{b}=0$ 的解）处能量密度和两个压强所满足的方程，用 $k$、$n$ 和 $N$ 表示。利用它们推导物态方程参数 $w^{(N)}=p^{(N)}/\rho$ 和 $w^{(n)}=p^{(n)}/\rho$ 在静态解处成立的表达式。

\textbf{2.}\quad 考虑 de Sitter 空间，采用度规取如下形式的坐标：
\begin{equation*}
\mathrm{d}s^{2} = -\mathrm{d}t^{2} + e^{2Ht}\left[\mathrm{d}x^{2}+\mathrm{d}y^{2}+\mathrm{d}z^{2}\right].
\tag{8.218}
\end{equation*}
对非共动观者（$x^{i}$ 不是常数）求解测地线方程，求出作为 $t$ 的函数的仿射参量。证明测地线在有限的仿射参量内到达 $t=-\infty$，从而证明这些坐标不能覆盖整个流形。

\textbf{3.}\quad 在附录 F 中我们讨论了 Raychaudhuri 方程。证明：把它应用到 Robertson--Walker 宇宙学上，Raychaudhuri 方程等价于第二弗里德曼方程 (8.68)。

\textbf{4.}\quad 考虑最佳拟合宇宙，其密度参数为 $\Omega_{\mathrm{R}0}=10^{-4}$、$\Omega_{\mathrm{M}0}=0.3$、$\Omega_{\Lambda0}=0.7$。在对数标度上画出三个 $\Omega_{i}$ 作为尺度因子 $a$ 的函数的图像，$a$ 从 $10^{-35}$ 取到 $10^{35}$。标出 Planck 时间、核合成和今天。

\textbf{5.}\quad 在平直时空中，物理尺寸固定的物体离得越远，所张的角就越小；在膨胀宇宙中未必如此。考虑红移 $z$ 处物理尺寸为 $L$ 的物体的角大小 $\theta(z)$。在物质主导的平直宇宙中，$\theta(z)/L$ 在什么红移处取极小？如果所有星系的跨度都至少有 $10~\mathrm{kpc}$（并且历来如此），这样一个宇宙中星系的最小角大小是多少？把你的结果既用 $H_{0}$ 表示，也代入 $H_{0}=70~\mathrm{km/s/Mpc}$ 给出数值。

\textbf{6.}\quad 在宇宙学中，我们倾向于把非相对论性粒子理想化为温度 $T$ 和压强 $p$ 都为零。实际上，随机运动会使它们具有一定的温度和压强，满足 $p\propto T\rho$。

\textbf{(a)}\quad 由有质量粒子组成的气体，其压强作为尺度因子的函数会如何衰减？

\textbf{(b)}\quad 设中微子具有质量 $m_{\nu}=0.1~\mathrm{eV}$，当前温度为 $T_{\nu 0}=2~\mathrm{K}$。中微子大约在什么红移从相对论性变为非相对论性？

\textbf{7.}\quad 设想宇宙在 Planck 时刻以均分状态（equipartition）出发（这样物质和辐射中的能量密度为 Planck 密度的量级，时间与空间的曲率半径为 Planck 长度的量级）。忽略任何空间不均匀性，计算一个正曲率宇宙能维持多久，以及一个负曲率宇宙在温度降到 $3~\mathrm{K}$ 时年龄多大。一个平直宇宙在温度降到 $3~\mathrm{K}$ 时年龄多大？一个平直宇宙到膨胀速率降至 $H_{0}=70~\mathrm{km\,s^{-1}\,Mpc^{-1}}$ 时年龄多大？
